\title{Riemann surfaces}
\author{\textit{Based on a lecture series by} Ruadhai Dervan.}
\date{Winter, 2023}
\maketitle

\marginpar{Lecture 1, 10/10/23} This course aims to be an introduction to
algebraic geometry, complex geometry, and parts of geometric analysis. We will
explain the theory of Riemann surfaces, these are smooth 2-manifolds with a lot
of structure which enables the notion of a holomorphic function (in one
variable.) We will use this setting to explore various facets of modern
geometry-- the study of differential geometry (differential forms and Hodge
structures), of algebraic geometry (compact Riemann surfaces are naturally
projective varieties), geometric analysis, and to a lesser extent topology and
complex analysis.

\subsection{About} This course on Riemann surfaces is a part of the SMSTC
training program for graduate students. The lectures were delivered in-person in
the University of Glasgow, and were transcribed live. The notes can be found
online at
\url{https://pas201.user.srcf.net/documents/2023-riemann-surfaces.pdf}. 

All errors and corrections should be communicated to by email to
\url{pas201@srcf.net}.

\subsection{Assessment} By presentation on some advanced topics after the
lectures. There will also be two exercise sheets, but these won't be marked.

\subsection{Prerequisites} We assume some familiarity with complex analysis,
basic notions of topology, (and not very advanced) differential geometry.

We recap some important ideas from complex analysis.

\begin{definition}
    A function \(f:U\to \mathbb{C}\) defined on an open \(U\subseteq
    \mathbb{C}\)  is \emph{holomorphic} if for all \(z_0\in U\), the limit
    \[f'(z_0)\coloneq \lim_{z\to z_0} \frac{f(z)-f(z_0)}{z-z_0}\]
    exists. Equivalently, writing \(z=x+iy\) and \(f(z)=u(x,y)+iv(x,y)\)  for
    real-valued functions \(u,v\), the function \(f\) is holomorphic if it is
    smooth and satisfies the Cauchy-Riemann equations
    \[
        \frac{\partial u}{\partial x}=\frac{\partial v}{\partial x}, \quad
        \frac{\partial u}{\partial y}=-\frac{\partial v}{\partial y}.
    \]
\end{definition}

\begin{theorem}
    A function \(f:U\to \mathbb{C}\) is holomorphic if and only if it is
    analytic, i.e.\ for all \(z_0\in U\) there is a neighbourhood
    \(U'\subseteq U\) of \(z_0\) such that \(f|_{U'}\) is given by a
    convergent power series \(f(z)=\sum_{n=0}^\infty c_n(z-z_0)^n\).
\end{theorem}

In particular, polynomials in \(z\) are holomorphic. The proof of the theorem
above uses the following important result.

\begin{theorem}[(Cauchy's integral formula)]
    Suppose \(f:U\to \mathbb{C}\) is holomorphic and write \(D=\{z\;|\;
    |z-z_0|\leq r\}\subset U\). Then for all \(a\) in the interior of \(D\), we
    have
    \[
        f(a)=\frac{1}{2\pi i}\int_{\partial D}\frac{f(z)}{z-a}dz
    \]
    where the integral is over the boundary of \(D\) (a circle oriented
    anticlockwise).
\end{theorem}


\begin{corollary}[(Identity principle)]
    Given two holomorphic functions \(f,g:U\to \mathbb{C}\) on an open connected
    domain \(U\subseteq {C}\), if we have \(f|_V=g|_V\) for a non-empty open set
    \(V\subseteq U\) then we have \(f=g\).
\end{corollary}

\section{Riemann surfaces}
The definition of a Riemann surface will be modelled after that of a smooth
manifold, except that we will use holomorphic transition functions instead.
\begin{definition}
    A Riemann surface is a Hausdorff topological space \(X\) with an open cover
    \(X=\bigcup_{\alpha}U_\alpha\), and for each \(\alpha\) a homeomorphism
    \(\phi_\alpha:U_\alpha \to \widetilde{U}_\alpha\) to an open subset
    \(\widetilde{U}_\alpha\subset \mathbb{C}\) such that whenever
    \(U_\alpha\cap U_\beta \neq \emptyset\), the function
    \(\phi_\alpha\circ \phi_\beta^{-1}\) is holomorphic where defined. We say
    each pair \((U_\alpha, \phi_\alpha)\) is a \textit{chart}, and the
    collection of all charts forms an \textit{atlas}.
\end{definition}

\begin{remark}
    A higher dimensional complex manifold is analogously defined, where we ask
    for the charts to go into \(\mathbb{C}^n\) with transition maps holomorphic
    in each coordinate.
\end{remark}

In practice, we work locally. For \(p\in X\), \(p\) is in some open
\(U_\alpha\) and then \(\phi_\alpha\) is a complex-valued function on
\(U_\alpha\) which we write as \(z\) . Thus locally we are doing complex
analysis. The transition functions being holomorphic ensures this local
coordinate is well-defined up to a holomorphic function, i.e.\ if another chart
provides a local coordinate \(w\) then \(w=f(z)\) for some holomorphic
invertible function \(f\).

\begin{definition}
    A function \(f:X\to \mathbb{C}\) is called \textit{holomorphic} if
    \(f|_{U_\alpha}\) is holomorphic (in the local coordinate) for all charts
    \(U_\alpha\subset X\).
\end{definition}

As opposed to differential geometry where a manifold has many smooth functions
on it, a Riemann surface might have no non-constant holomorphic functions (eg.\
if it is compact). On the other hand, maps between Riemann surfaces are usually
abundant.

\begin{definition}
    Let \(X\) be a Riemann surface with atlas \(\{(U_\alpha, \phi_\alpha)\}\)
    and \(Y\) a Riemann surface with atlas \(\{(V_\alpha,\phi_\alpha)\}\). A
    map \(f:X\to Y\) is \textit{holomorphic} if for all pairs of indices
    \(\alpha, \beta\), the map \(\phi_\beta\circ f\circ \phi_\alpha^{-1}\) is
    holomorphic where defined.
\end{definition}

We say a bijection \(f\) is a biholomorphism if \(f^{-1}\) is holomorphic, and
we usually treat biholomorphic Riemann surfaces as identical.

\begin{example} The following constructions are standard.
    \begin{enumerate}
        \item The complex plane \(\mathbb{C}\) itself is a complex manifold,
            with a single chart (with the identity map). Likewise any open
            \(U\subset \mathbb{C}\) is a complex manifold with a single chart.

        \item  The \(2\)-sphere \(S^2=\mathbb{C}\cup \{\infty\}\) is a complex
            manifold called the \textit{Riemann sphere}, with two charts
            \(S^2\setminus \{\infty\} \cong \mathbb{C}\) (via the
            identity) and \(S^2\setminus \{0\}\cong \mathbb{C}\)
            (via \(z\mapsto \frac{1}{z}\)).

        \item If \(\pi:S\to X\) is a topological covering space and \(X\) is a
            Riemann surface, then \(S\) canonically obtains the structure of a
            Riemann surface that makes \(\pi\) holomorphic.

        \item Conversely we can construct Riemann surfaces as quotients--
            consider a lattice \(\Lambda \subset \mathbb{C}\) (i.e.\ a rank two
            additive subgroup generated by two \(\mathbb{R}\)-linearly
            independent complex numbers). We can take the quotient topological
            space \(\mathbb{C}/\Lambda\) which is homeomorphic to a torus, and
            equip it with charts on which each open set is small enough to lie
            within a fundamental domain of the \(\Lambda\)-action after pulling
            back to \(\mathbb{C}\) . This is an \textit{elliptic curve}.
    \end{enumerate}
\end{example}

We will now discuss two important classes of examples in detail.

\subsection{Affine algebraic curves}
Take a polynomial \(p(z,w):\mathbb{C}^2\to \mathbb{C}\) in two complex
variables, and consider the set \(S=\{(z,w)\in \mathbb{C}^2\;|\;
p(z,w)=0\}\). This is an \textit{algebraic curve}. When is this a Riemann
surface? This has to do with whether the space is singular or not-- in
particular we want to be able to locally write \(w\) as a function of \(z\) (or vice
versa). The implicit function theorem deals with precisely this.

\begin{theorem}[(Implicit function theorem)]
    For \(p,S\) as above and \((z_0,w_0)\in S\), if \(\frac{\partial
    p}{\partial w}(z_0,w_0)\neq 0\) then there are open balls \(B_1\ni z_0\) and
    \(B_2\ni w_0\) in \(\mathbb{C}\) such that \(S\cap (B_1\times
    B_2)\subseteq \mathbb{C}^2\) is the graph of a holomorphic function
    \(\sigma: B_1\to B_2\), i.e.\ \(S\cap (B_1\times B_2) = \{(z,\sigma
    z)\;|\; z\in B_1\}.\)
    In particular, \((z_0,w_0)\) has a neighbourhood in \(S\) that is
    homoeomorphic to \(B_1\).
    \begin{proof}
        We may assume \((z_0,w_0)=(0,0)\). First consider the holomorphic
        function \(p(0,w)\), this is a polynomial in one variable with
        \(p(0,w)=0\) for \(w=0\) and \(\frac{\partial p}{\partial w}(0,0)\neq 0\)
        by hypothesis. Thus \(0\) is a simple zero. As zeros of holomorphic
        functions are isolated, we can find a \(\delta_2 > 0\) with
        \(p(0,w)\neq 0\) whenever \(0<|w|<2\delta_2\). By continuity of \(p\),
        there is a \(\delta_1>0\) such that \(p(z,w)\neq 0\) whenever
        \(|w|=\delta_2\), \(|z|\leq \delta_1\). Let \(B_1,B_2\) be the open
        balls around the origin of radii \(\delta_1,\delta_2\) respectively.
        Then for \((z,w)\in S\cap (B_1\times B_2)\), will write \(w\) as a
        function of \(z\).

        For this, we use that for holomorphic \(f\) and a simple closed
        curve \(\gamma\), we have
        \[
            \frac{1}{2\pi i}\int_\gamma \frac{f'(w)}{f(w)}dw =
            \#\{\text{zeros of \(f\) in the region bounded by \(\gamma\).}\}
        \]
        Consider the holomorphic functions \(f_z(w)=p(w,z)\) for \(z\in B_1\),
        and the contour \(\gamma(t)=\delta_2\cdot e^{2\pi i t}\).
        Then \(N(z)=\frac{1}{2\pi i}\int_\gamma \frac{f_z'(w)}{f_z'(w)}dw\) is
        an integer that varies continuously with \(z\), and so considering
        \(z=0\) gives us \(N(z)=1\) identically. Thus for each \(z\in B_1\)
        there is a unique \(w\in B_2\) with \(f_z(w)=0\). In fact the zero is
        then given by
        \[\sigma(z)=\frac{1}{2\pi i}\int_\gamma w\frac{f'_z(w)}{f_z(w)}dw,\]
        which expresses \(w\) as a holomorphic function of \(z\) as required.
    \end{proof}
\end{theorem}

\begin{corollary}
    For \(p,S\) as above, \(S\subset \mathbb{C}^2\) is a Riemann surface if
    for all \((z_0,w_0)\in S\), either \(\frac{\partial p}{\partial
    z}(z_0,w_0)\neq 0\) or \(\frac{\partial p}{\partial
    w}(z_0,w_0)\neq 0.\) In this case, we say \(S\) is the \emph{affine
    algebraic curve} defined by \(p\).
\end{corollary}


\subsection{Projective curves}
Recall the projective plane is \(\mathbb{P}^2=\mathbb{C}^3\setminus \{0\}/\sim\)
where \((z_1,z_2,z_3)\sim (\lambda z_2,\lambda z_2, \lambda z_3)\) for all
\(\lambda\in \mathbb{C}^\times\). This is compact-- this fact is easily seen by
observing that \(\mathbb{P}^2\) is the quotient of the (compact) subset of unit
vectors in \(\mathbb{C}^3\) by the action of the compact group \(S^1\subset
\mathbb{C}^\times\). Points in \(\mathbb{P}^2\) can be labelled by the
homogeneous coordinates \([z_0:z_1:z_2]\), likewise for higher dimensional
projective spaces.

Write \(U_0=\{[z_0:z_1:z_2]\;|\; z_0\neq 0\}\) and note that this is
homeomorphic to \(\mathbb{C}^2\) since all points have homogeneous coordinates
of the form \([1:z:w]\). The complement \(\mathbb{P}^2\setminus U_0\) is
homeomorphic to \(\mathbb{P}^1\), so we can say \(\mathbb{P}^2=\mathbb{C}^2\cup
\mathbb{C}\cup  \{\infty\}\).

Note polynomials in \(z_0,z_1,z_2\) don't take well-defined values on
\(\mathbb{P}^2\), but the vanishing locus of a homogeneous polynomial is
nonetheless well-defined. Indeed, if \(p(z_0,z_1,z_2)\) is homogeneous of degree
\(d\) then for \(\lambda\in \mathbb{C}^\times\) we have \(p(\lambda z_0,
\lambda z_1, \lambda z_2) = \lambda^dp(z_0,z_1,z_2)\). Thus it makes sense to
ask when a homogeneous polynomial \(p\) defines a Riemann surface inside
\(\mathbb{P}^2\).

\begin{proposition}
    If \(p(z_0,z_1,z_2)\) is a homogeneous polynomial such that
    the only solution to \(\frac{\partial p}{\partial z_0}=\frac{\partial
    p}{\partial z_1}=\frac{\partial p}{\partial z_2}= 0\) is \((0,0,0)\), then
    the vanishing set \(S\) of \(p\) inside \(\mathbb{P}^2\) is naturally a
    compact Riemann surface.

    \begin{proof}
        We will show that \(S\cap U_0 \subset \mathbb{C}^2\) is a Riemann
        surface by using the implicit function theorem. It will be true
        similarly and in a compatible way for \(U_1\) and \(U_2\), giving the
        required charts. The compactness then follows since \(\mathbb{P}^2\) is
        compact and \(S\subset \mathbb{P}^2\) is closed.

        To use the implicit function theorem, we need to rule out the
        possibility that there is a \([1:z:w]\in S\) where \(\frac{\partial
        p}{\partial z_1}(1,z,w)=\frac{\partial p}{\partial z_2}(1,z,w)=0\). Note since \(p\) is homogeneous, say of degree \(d\), we have Euler's
        identity \(\sum_{i=0}^2 z_i \frac{\partial p}{\partial z_i} = d\cdot
        p\).  Thus if \(\frac{\partial p}{\partial z_1}(1,z,w)=\frac{\partial
        p}{\partial z_2}(1,z,w)=0\)  we must have \(\frac{\partial p}{\partial
        z_0}(1,z,w)=0\), contradicting the hypothesis that \((0,0,0)\) is the
        only point where this happens.
     \end{proof}
\end{proposition}

\begin{example}
    Let \(p(z,w)=z^5+w^5-zw-1\). We homogenize to get
    \(p(z_0,z_1,z_2)=z_1^5+z_2^5-z_0^3z_1z_2-z_0^3\), with vanishing set
    \(S\subset \mathbb{P}^2\). Then \(S\cap U_0\) is the vanishing locus of the
    original polynomial, while \(S\cap (\mathbb{P}^2\setminus U_0)\) has five
    points \([0:1:e^{\frac{-2\pi i k}{5}}]\), \(k=0,\ldots,4\).
\end{example}

\begin{remark}
    One can show that a Riemann surface is always orientable, i.e.\ there is an
    atlas where all the Jacobians of the transition functions have positive
    determinant. Thus compact Riemann surfaces are given by adding a holomorphic
    structure to an \(n\)-holed torus, \(n\geq 0\). Donaldson sketches a proof
    of this using morse theory.
\end{remark}

\begin{remark}
    Historically Riemann surfaces arose as the natural domains of `germs' of
    multi-valued holomorphic functions. We won't go into this.

\end{remark}


\section{Holomorphic maps}
Holomorphic functions are well understood, so we can be ambitious and aim to
have a structure theory of maps between Riemann surfaces (which are locally
holomorphic). This can shed light on important structure of the Riemann
surfaces-- for instance a compact Riemann surface has no non-constant
holomorphic maps to \(\mathbb{C}\). By contrast, affine curves \(S\subset
\mathbb{C}^2\) have many holomorphic functions.

\begin{example}[Maps between torii]
    Consider \(\mathbb{C}/\Lambda\) for \(\Lambda\) a lattice. When does a
    linear map \(G:\mathbb{C}\to \mathbb{C}\) induce a holomorphic map
    \(\mathbb{C}/\Lambda \to \mathbb{C}/\Lambda\)?
    \marginpar{Lecture 2, 17/10/23}Clearly we need
    \([G(z)]=[G(z+\lambda)]\in \mathbb{C}/\Lambda\) for all \(\lambda \in
    \Lambda\), so if \(G(z)=az+b\) then this amounts to the requirement that
    \(a\Lambda \subset \Lambda\). Thus for instance \(a=1\) works, i.e.\
    translations always induce maps \(\mathbb{C}/\Lambda \to
    \mathbb{C}/\Lambda.\) For a non-trivial example, consider \(\Lambda=\langle
    1,i\rangle\) and \(G(z)=(1+i)z\). This gives a \(2:1\) cover from the torus
    to itself.
\end{example}

\begin{example}
    Just like holomorphic functions, we have a notion of \textit{meromorphic
    functions} on a Riemann surface i.e.\ \(f:X\to \mathbb{C}\cup \{\infty\}\)
    that are not identically \(\infty\) such that for all charts \(\psi:U\to
    \psi(U)\subseteq \mathbb{C}\),  the function \(f\circ \psi^{-1}\) is
    meromorphic on \(\psi(U)\). Viewing \(\mathbb{C}\cup \infty\) as the Riemann
    sphere, we see that meromorphic functions are just
    holomorphic maps to \(\mathbb{P}^1\).
\end{example}

To get general theory, we use the inverse function theorem.

\begin{theorem}[(Inverse function theorem)]
    Let \(U\ni 0\) be an open subset of \(\mathbb{C}\), and \(f:U\to
    \mathbb{C}\) holomorphic such that \(f'(0)\neq 0\). Then there is a possibly
    smaller neighbourhood \(U'\subseteq U\) of \(0\) such that \(f:U'\to
    f(U')\) has a holomorphic inverse.

    \begin{proof}
        Similar to implicit function theorem.
    \end{proof}
\end{theorem}

\begin{corollary}
    Suppose \(U\ni 0\) is an open subset of \(\mathbb{C}\) and \(f:U\to
    \mathbb{C}\) is a holomorphic function that vanishes on \(0\) but is not
    identically zero. Then there is some \(k\geq 1\) such that on a possibly
    smaller neighbourhood \(U'\subseteq U\) of \(0\), we have \(f(z)=g(z)^k\)
    for a holomorphic function \(g:U'\to \mathbb{C}\) with \(g'(0)\neq 0\).

    \begin{proof}
        This is automatic-- since \(f\) is analytic and vanishes at the origin,
        we have \(f(z)=z^k(a_0+a_1z+\ldots)\) for \(k\geq 1\), \(a_0\neq 0\).
        Thus \(f(z)=z^kh(z)\) for some holomorphic function \(h\) with
        \(h(0)\neq 0\). Then by the inverse function theorem applied to
        \(z\mapsto z^k\), the function \(z\mapsto z^{1/k}\) is well-defined on a
        neighbourhood of \(h(0)\) and we can set \(g(z)=z
        \cdot h(z)^{1/k}\) to get the required result.
    \end{proof}
\end{corollary}

\begin{proposition}
    Suppose \(F:X\to Y\) is a non-constant holomorphic map between connected
    Riemann surfaces. Then for all \(x\in X\), there exist neighbourhoods \(U\ni
    x\), \(V\ni f(x)\) and biholomorphisms \(\phi:U\to \phi(U)\subseteq
    \mathbb{C}\), \(\psi:U\to \psi(U)\subseteq \mathbb{C}\) such that
    \(\phi(x)=0\) and \(\psi\circ F\circ \phi^{-1}\) is
    given by \(z\mapsto z^k\) for some integer \(k\geq 1\).

    \begin{proof}
        This is a consequence of the above results. Start with any charts
        \(\phi', \psi'\) around \(x, f(x)\) respectively and set \(f=\psi'\circ
        F\circ \phi'^{-1}\). This is a holomorphic map, and after composing the
        charts with translations we may assume \(f(0)=0\). Then \(f(z)=g(z)^k\)
        for some holomorphic \(g\) satisfying \(g(0)=0\), \(g'(0)\neq 0\).
        Define \(\phi=g\circ \phi'\), and note this is holomorphic and
        invertible on a possibly smaller open neighbourhood of \(x\). Thus we
        have a new chart in which \(\psi'\circ F\circ \phi^{-1}\) has the
        required form.
    \end{proof}
\end{proposition}

Note the \(k\) in the above proposition depends on both \(F\) and \(x\), but is
independent of the choice of charts. To emphasise the dependence on \(x\) we
will often write \(k_x\).

\begin{definition}
    In the situation of the theorem above, we say \(x\in X\) is a
    \textit{ramification point} of \(F\) if \(k_x>1\). The subset of all
    ramification points \(R\subseteq X\) is called the \textit{ramification
    locus}, and its image \(F(R)\subseteq Y\) is called the \textit{branch
    locus}.
\end{definition}

Since the ramification locus \(R\) is defined locally by \(x\in R\iff F'(x)=0\)
and since \(F'\) is holomorphic in each chart, we see that \(R\) is discrete
(and hence finite if \(X\) is compact). Likewise if \(Y\) is compact (or \(F\)
is proper), then the branch locus \(F(R)\) is finite (resp.\ discrete).

Now assume \(F\) is proper, and consider the preimage \(F^{-1}(y)\subset X\) of
a point \(y\in Y\). This set is discrete and compact, hence is finite. So we can
count it. Define (temporarily) the \textit{degree of \(F\) at \(y\)} to be
\[
    d(y)=\sum\nolimits_{x\in F^{-1}(y)}k_x.
\]

This is a temporary definition because in fact it does not depend on \(y\) at
all. Thus we define the \textit{degree of \(F\)} to be its degree at any \(y\in
Y\).

\begin{proposition}
    The integer \(d(y)\) does not depend on \(y\).

    \begin{proof}
        We will show the function \(d:Y\to \mathbb{Z}\) is locally constant on
        \(Y\), hence constant since \(Y\) is connected. Given \(y\in Y\) with
        preimages \(x_1,\ldots,x_n\in X\), pick neighbourhoods \(V\ni Y\),
        \(U_i\ni x_i\) sufficiently small so that \(U_i\cap U_j=\emptyset\) for
        \(i\neq j\), \(F(U_i)\subseteq V\), and the
        restriction \(F:U_i\to V\) is given by \(z\mapsto z^{k_i}\) in local
        coordinates. Here \(k_i=k_{x_i}\), and \(i=1,\ldots,n\). We wish to find a
        possibly smaller neighbourhood \(W\ni y\) on which \(d\) is constant.

        Claim that there is a neighbourhood \(W\ni y\) such that
        \(F^{-1}(W)\subseteq U_1\cup U_2\cup \ldots \cup U_n\). The properness
        hypothesis is essential here. Indeed, the natural map \(\mathbb{C}\sqcup
        \mathbb{P}^1\to \mathbb{P}^1\) is not proper and so while \(\infty \in
        \mathbb{P}^1\) has one pre-image, any point in the neighbourhood of
        \(\infty\) has two. Now having seen why the claim is non-trivial, we
        will prove it using point-set topology arguments. Choose a sequence of
        open sets \(V_1\supseteq V_2 \supseteq \ldots \ni y\) such that \(\bigcap
        \overline{V_i}=\{y\}\). If the claim were false, then in particular
        \(W=V_i\) does not suffice and we can choose a point
        \(z_i\in F^{-1}(V_i)\setminus (U_1\cup  \ldots \cup U_n)\) for each
        \(V_i\). This gives a sequence \(z_1,z_2,\ldots\) in the set
        \(F^{-1}(\overline{V_1})\setminus (U_1\cup \ldots \cup U_n)\), which by the
        properness hypothesis is compact. Thus there is a subsequence converging
        to some \(z\notin U_1\cup \ldots \cup U_n\). But by continuity of \(F\) and
        choice of \(V_i\) , we have \(F(z)=y\). This is a contradiction since we
        know all pre-images of \(y\) are contained in some \(U_i\).

        Having proved the claim, the result follows since for any \(y'\in W\),
        we have
        \[
            d(y')=\sum\nolimits_{x\in F^{-1}y'}k_x = \sum\nolimits_{i=1}^n
            \sum\nolimits_{x\in
            F^{-1}y'\cap U_i} k_x
        \]
        but locally on \(U_i\), the map is \(z\mapsto z^{k_i}\) and hence we
        have \(\sum\nolimits_{x\in F^{-1}y'\cap U_i} k_x =k_i\) since either
        \(y'=y\) (and it has one preimage in \(U_i\)) or \(y'\neq y\) (and it
        has \(k_i\) simple preimages).
    \end{proof}
\end{proposition}

\begin{corollary}
    Any proper non-constant holomorphic map between Riemann surfaces is
    surjective.
\end{corollary}

\begin{remark}
    The above corollary could also have been obtained from the open mapping
    principle.
\end{remark}

We next use this to obtain topological consequences-- in particular, counting
poles of meromorphic functions leads to a classification of compact Riemann
surfaces.

\begin{corollary}
    Let \(X\) be a compact Riemann surface. If there is a meromorphic function
    \(f:X\to \mathbb{C}\cup \{\infty\}\) with a single simple pole, then \(f\)
    provides a biholomorphism \(X\cong \mathbb{P}^1\).

    \begin{proof}
        Note we can view \(f\) as a holomorphic map \(X\to \mathbb{P}^1\). Since
        \(X\) is compact, \(f\) is proper and hence has degree
        \(d=d(\infty)=1\). Thus \(F\) is a continuous bijection between compact
        Hausdorff spaces, so by the topological inverse function theorem, \(f\)
        is a homeomorphism. To show its a biholomorphism it suffices to show the
        inverse continuous map \(f^{-1}\) is in fact holomorphic. But this is
        clear since in a local chart around any \(x\in X\), \(f\) has
        non-vanishing derivative and so the local inverse is holomorphic by the
        inverse function theorem.
    \end{proof}
\end{corollary}

This concludes for now the study of maps between Riemann surfaces. In higher
dimensions, there are many more interesting maps but codimension one phenomena
are nonetheless controlled by the same rules which describe maps between Riemann
surfaces.

\section{Calculus on surfaces}
We wish to eventually develop a theory of holomorphic differential forms and de
Rham cohomology on Riemann surfaces. Before that, we give an account of the
differential geometry of smooth surfaces.

\subsection{Tangent and cotangent spaces}
There are various ways to define vector fields and differential forms on
surfaces, for instance these can be realised most naturally as sections of the
tangent and cotangent bundles respectively. Here we will take a more
abstract approach, defining tangent vectors as infinitesimal paths and
differential forms as derivations of smooth functions.

\begin{lemma}
    Let \(U\subseteq \mathbb{R}^2\) be an open neighbourhood of \(0\),
    \(f:U\to\mathbb{R}\) a smooth function, and
    \(\gamma_1,\gamma_2:\mathbb{R}\to U\) two smooth paths with
    \(\gamma_1(0)=\gamma_2(0)=0\). Let \(\chi:U\to V\subseteq \mathbb{R}^2\) be
    a diffeomorphism with \(\chi(0)=0\), and set \(\tilde{f}=f\circ \chi^{-1}\),
    \(\tilde{\gamma}_1=\chi\circ \gamma_1\), \(\tilde{\gamma}_2=\chi\circ
    \gamma_2\). Then the following hold.
    \begin{enumerate}[(i)]
        \item If both partial derivatives \(\frac{\partial f}{\partial x_1}\)
        and \(\frac{\partial f}{\partial x_2}\) vanish at \(0\) then the same
        is true for \(\tilde{f}\).
        \item If the derivatives \(\frac{d\gamma_1}{dt}(0)\) and
        \(\frac{d\gamma_2}{dt}(0)\) agree then the same is true for
        \(\tilde{\gamma}_1\) and \(\tilde{\gamma}_2\).
    \end{enumerate}
    \begin{proof}
        These are straightforward consequences of the chain rule.
    \end{proof}
\end{lemma}

Given the diffeomorphism invariance of these notions, we can make the following
definitions. Let \(S\) be a (real, smooth) surface with a point \(p\in S\),
\(f\)  a smooth real-valued function defined on a neighbourhood of \(p\), and
\(\gamma_1,\gamma_2: \mathbb{R}\to S\) smooth paths with
\(\gamma_1(0)=\gamma_2(0)=p\).

\begin{definition}
    We say \(f\) is \textit{constant to first order} at \(p\) if \(\frac{\partial
    f}{\partial x_1}\) and \(\frac{\partial f}{\partial x_2}\) both vanish at
    \(p\) in some local chart. Likewise we say \(\gamma_1,\gamma_2\) are
    \textit{equal to first order} at \(p\) if their derivatives agree in a
    local chart.
\end{definition}

\begin{definition}
    The \textit{cotangent space} \(T^\ast_p S\)  of \(S\) at \(p\) is the set of equivalence
    classes of smooth functions on open neighbourhoods of \(p\in S\), where
    \(f_1\) and \(f_2\) are equivalent if and only if \(f_1-f_2\) is constant
    to first order at \(p\).

    Dually, the \textit{tangent space} \(T_pS\)  of \(S\) at \(p\) is the set of
    equivalence classes of smooth paths \(\gamma:\mathbb{R}\to S\) with
    \(\gamma(0)=p\) where two paths are equivalent if they agree to first
    order at \(p\).
\end{definition}

Note there is an obvious map \(C^\infty(S)\to T_p^\ast S\) sending a function
\(f\) to its equivalence class, for which we will write \([df]_p\) (and omit
the brackets whenever convenient). This gives \(T_p^\ast S\) a natural vector
space structure, and in fact we can check that it is two dimensional. Indeed
choosing local coordinates \(x_1,x_2\) around \(p\), these provide smooth
functions and so we have two classes \([dx_1]_p, [dx_2]_p\in T_p^\ast S\). One checks that
\[
    [df]_p = \frac{\partial f}{\partial x_1}(p)\cdot [dx_1]_p + \frac{\partial
    f}{\partial x_2}(p)\cdot [dx_2]_p,
\]
hence any choice of local coordinates provides a basis for the cotangent space
at the point. Similarly \(T_pS\) is a vector space, by choosing linear representatives of paths.

The two vector spaces are naturally dual to each other-- if
\(\gamma:\mathbb{R}\to S\) is a path with \(\gamma(0)=p\) and \(f\) is a smooth
function defined on a neighbourhood of \(p\), then the derivative \((f\circ
\gamma)'(0)\in \mathbb{R}\) depends only on the classes \([df]_p\in T_p^\ast
S\), \([\gamma]\in T_pS\). Thus we have a perfect bilinear pairing
\begin{align*}
    T_pS\times T_p^\ast S \to \mathbb{R}, \qquad
    ([\gamma], [f]) \mapsto (f\circ\gamma)'(0)
\end{align*}
giving an isomorphism \(T_p^\ast S=\Hom_\mathbb{R}(T_pS, \mathbb{R})\).

As a set, the \textit{cotangent bundle} of \(S\) is given by \(T^\ast
S=\bigsqcup_p T^\ast_pS\). This can naturally be given the structure of a vector
bundle over \(S\), but we won't go into the details of this. Smooth sections of
this bundle are called \(1\)-forms, but since we did not give the vector
bundle structure on \(T^\ast S\) we will define what it means to be a
\(1\)-form explicitly.

\begin{definition}
    A smooth \(1\)-form \(\alpha\) on \(S\) is a map \(\alpha:S\to T^\ast S\)
    satisfying \(\alpha(p)\in T_p^\ast S\) for all \(p\), and varying smoothly
    with \(p\) in the following sense: in local coordinates \((x_1,x_2)\)
    around any \(p\in S\), we can write \(\alpha=\alpha_1 dx_1 + \alpha_2 dx_2\)
    for smooth functions \(\alpha_1, \alpha_2\).
\end{definition}

This notion of smoothness is well-defined. Indeed if \((y_1,y_2)\) were another
pair of local coordinates around \(p\in S\), then the corresponding cotangent
vectors are related as
\[
    dx_i  = \frac{\partial x_i}{\partial y_1}dy_1 + \frac{\partial x_i}{\partial
    y_2} dy_2
\]
wherever defined so the coefficients of the \(1\)-form \(\alpha\) remain smooth.

(Co)tangent vectors behave functorially with respect to smooth maps between
surfaces-- a map \(F:S\to Q\) naturally induces linear maps
\begin{align*}
    dF_p: T_pS \to T_{F(p)}Q, \quad [\gamma] \mapsto [F\circ \gamma] \\
    F_p^\ast:T^\ast_{F(p)}Q \to T^\ast_pS, \quad [f]\mapsto [f\circ F].
\end{align*}
Note the behaviour is covariant on tangent vectors and contravariant on
cotangent vectors. This pullback map extends to \(1\)-forms naturally, so if
\(\alpha\) is a \(1\)-form on \(Q\) then the \(1\)-form \(F_p^\ast \alpha\) given
by \((F_p^\ast\alpha)(p)= F_p^\ast(\alpha(F(p)))\) is still smooth.

\subsection{Integration along curves}\marginpar{Lecture 3, 24/10/23}
One of the things that makes differential forms useful is that they can be
integrated along curves. In the simplest case, suppose \(\gamma:[0,1]\to
S\) is a smooth curve in our surface whose image lies in a chart with local
coordinates \(x_1,x_2\) and suppose \(\alpha\) is a
\(1\)-form given on this chart as \(\alpha_1dx_1+\alpha_2dx_2\). Writing
\(\gamma(t)=(\gamma_1(t),\gamma_2(t))\) in these local coordinates, we can
define the integral of \(\alpha\) along \(\gamma\) as
\[
    \int_\gamma \alpha = \int_0^1(\alpha_1\cdot
    \gamma_1'(t)+\alpha_2\cdot \gamma_2'(t))dt.
\]
One checks that this is independent of local coordinates chosen. For a general
curve \(\gamma\), we cut it into smaller curves \(\gamma_1,\ldots,\gamma_n\) each
of which is contained in a single coordinate chart, and then define
\(\int_\gamma \alpha = \sum_{i=1}^n \int_{\gamma_i}\alpha\).

The fundamental theorem of calculus holds i.e.\ \(\int_\gamma df =
f(1)-f(0)\) for any smooth function \(f\). This can be proven locally on charts.

\subsection{Higher differential forms} If we wish to analogously extend the
theory of integration to surfaces, we need to define \(2\)-forms. Just as a
function gives a \(1\)-form, there is an operator taking a \(1\)-form to a
\(2\)-form.  But this is more involved; the classical motivation for the
definition is as follows.

Given a \(1\)-form \(\alpha\) on a surface \(S\), we
ask when is \(\alpha=df\) for some function \(f\)? When \(S=\mathbb{R}^2\), we
have \(\alpha=\alpha_1dx_1 +\alpha_2dx_2\) for a choice of coordinates. If
\(\alpha=df\) then
\(
    \frac{\partial f}{\partial x_1}
    = \alpha_1, \quad  \frac{\partial f}{\partial x_2} = \alpha_2.
\)
Thus by the symmetry of mixed partial derivatives, a necessary
condition for \(\alpha\) to be of the form \(df\) is that the
function \(R=\frac{\partial \alpha_1}{\partial x_2} - \frac{\partial
\alpha_2}{\partial x_1}\) vanishes identically. This is also sufficient-- if \(R=0\) then we can find \(f\) as follows.
Define functions \(f_1,f_2\) by
\begin{align*}
    f_2(x_1,x_2)=\int_0^{x_1}\alpha_1(t,0)+\int_0^{x_2}\alpha_2(x_1,t)dt,\qquad
    f_1(x_1,x_2)=\int_0^{x_2}\alpha_1(0,t)+\int_0^{x_1}\alpha_2(t,x_2)dt.
\end{align*}
By construction, \(\frac{\partial f_i}{\partial x_i}=\alpha_i\). Then Green's
theorem applied to the rectangle \(V\) with vertices \((0,0)\), \((x_1,0)\),
\((0,x_2)\), \((x_1,x_2)\) shows \(f_1(x_1,x_2)-f_2(x_1,x_2)=\int_V
Rdx_1dx_2\). Thus if \(R\) vanishes, then \(f_1=f_2\) and this is the required
function \(f\).

We generalise this to arbitrary surfaces, with the motivation that the problem
in question is controlled by some anti-symmetric function. With this in mind,
recall that for a real vector space \(E\) we define \(\Exter^2E^\ast\) to be the
space of skew-symmetric bilinear maps \(E\times E\to \mathbb{R}\). There is a
natural wedge-product operator
\begin{align*}
    \wedge:E^\ast \times E^\ast &\to \Exter^2 E^\ast \\
    (\alpha, \beta)(e,f) &= \alpha(e)\beta(f)-\alpha(f)\beta(e).
\end{align*}
It is clear that this product is linear in each variable and \(\alpha\wedge
\beta = -\beta\wedge \alpha\). If \(E\) is \(2\)-dimensional, one checks that
\(\Exter^2E^\ast\) is \(1\)-dimensional and any basis \(\alpha_1,\alpha_2\) of
\(E^\ast\) induces a basis \(\alpha_1\wedge \alpha_2\) of \(\Exter^2E^\ast\).

We are interested in the case \(E=T_pS\) for a smooth surface \(S\), so
\(E^\ast\) is the cotangent space \(T_p^\ast S\). A choice of local
coordinates \(x_1,x_2\) induces a basis element \(dx_1\wedge dx_2\) for
\(\Exter^2T_p^\ast S\).

\begin{definition}
    A (smooth) \(2\)-form \(\rho\) on \(S\) is a map from \(S\) to
    \(\bigcup_{p\in S}\Exter^2T_p^\ast S\) such that \(\rho(p)\in
    \Exter^2T_p^\ast S\) for all \(p\in S\), and such that \(\rho\) varies
    smoothly in the sense that locally we have \(\rho=R(x_1,x_2)dx_1\wedge
    dx_2\) for some smooth function \(R\).
\end{definition}

This is well defined independent of chart; for a different choice of coordinates
\(y_1,y_2\), \(\rho\) has the form \( \rho=R(y_1,y_2)\cdot J(y_1,y_2) \cdot
dy_1\wedge dy_2 \) where \(J(y_1,y_2)\) denotes the Jacobian of \(y_1,y_2\) with
respect to \(x_1,x_2\), and \(dy_i=\frac{\partial y_i}{\partial x_1}dx_1 +
\frac{\partial y_i}{\partial x_2}dx_2\).

One can pull back \(2\)-forms just like one does for \(1\)-forms. In particular
if \(F:S\to Q\) is a smooth map and \(\rho\) is a \(2\)-form on \(Q\) then one
obtains a \(2\)-form \(F^\ast \rho\) on \(S\), expressed in local coordinates in
a very similar way.

\subsection{The exterior derivative} We write \(\Omega^i(S)\) for the set of smooth \(i\)-forms on \(S\), where a
smooth \(0\)-form is nothing but a function. We have already seen that there is
a differential \(d:\Omega^0(S)\to \Omega^1(S)\) sending a function \(f\) to
\(df\). The content of the following result is that this extends nicely.

\begin{lemma}
    There is a unique way to define an \(\mathbb{R}\)-linear map
    \(d:\Omega^1(S)\to \Omega^2(S)\) such that
    \begin{enumerate}[(i)]
        \item For \(\alpha\in \Omega^1(S)\) and open set \(U\subset S\), we have
        \(\alpha|_U\in \Omega^1(U)\) and \((d\alpha)|_U=d(\alpha|_U)\).
        \item If \(\alpha_1=\alpha_2\) on an open set \(U\subset S\), then
        \(d\alpha_1=d\alpha_2\) on \(U\).
        \item For any function \(f\in \Omega^0(S)\), we have \(d(df)=0\).
        \item If \(f\in \Omega^0(S)\) and \(\alpha\in \Omega^1(S)\), we have
        \(d(f\cdot \alpha)=df\wedge \alpha + f\cdot d\alpha\).
    \end{enumerate}

    \begin{proof}
        If we have have defined \(d\) with these properties, then for a
        \(1\)-form \(\alpha=\alpha_1\cdot dx_1+\alpha_2\cdot dx_2\) in local
        coordinates on an open \(U\subset S\), we have
        \begin{align*}
            d\alpha &= d(\alpha_1\cdot dx_1 +\alpha_2 \cdot dx_2) \\
                    &= d\alpha_1 \wedge dx_1 + d\alpha_2 \wedge dx_2 \\
                    &= \ldots \\
                    &= \left(\frac{\partial \alpha_2}{\partial x_1} -
                    \frac{\partial \alpha_1}{\partial x_2}\right) dx_1\wedge
                    dx_2.
        \end{align*}
        This shows uniqueness, and also gives a way to define \(d\) after one
        checks it is independent of charts.
    \end{proof}
\end{lemma}

Having expressed \(2\)-forms as derivatives, we now see how to integrate them.
For this we need to assume \(S\) is oriented (so there are charts such that the
coordinate-change Jacobians are everywhere positive). Let \(\rho\in
\Omega^2(S)\) have compact support in a single chart. Then we can write
\(\rho=R(x_1,x_2)dx_1\wedge dx_2\) in local coordinates, and define
\[
    \int_S\rho = \int_{\mathbb{R}^2}R(x_1,x_2)\cdot dx_1 dx_2
\]
where the right hand side is a Lebesgue integral. If \(y_1,y_2\) is another
oriented chart with positive coordinate-change Jacobian, then that this is
well-defined is the usual transformation law for multiple integrals.

To integrate in general, we need the following technical lemma which says
partitions of unity exist.

\begin{lemma}
    Let \(K\subset S\) be compact and let \(U_1,\ldots,U_n\) be open sets with
    \(K\subseteq U_1\cup \ldots \cup U_n\). Then there are smooth non-negative
    functions \(f_1,\ldots,f_n\) on \(S\) with \(f_1+\ldots+f_n=1\) on \(K\),
    such that each \(f_i\) is supported on \(U_i\).
\end{lemma}

\begin{definition}
    For a \(2\)-form \(\rho\in \Omega^2(S)\) with compact support, we define
    \(\int_S\rho = \sum_i\int_S(f_i\cdot \rho)\) where the \(f_i\) form a
    partition of unity over \(\text{Supp}(\rho)\) and each \(f_i\) is supported
    on a coordinate chart.
\end{definition}

To see this makes sense, since \(\sum_i f_i =1\) we have \(\rho=\sum_i
(f_i\cdot \rho)\) and each \(f_i\cdot \rho\) is supported on a coordinate chart
so can be integrated locally. The integral is independent of the choice of
partition of unity by linearity of Lebesgue integrals.

\begin{theorem}[(Stokes' theorem)]
    If \(\alpha\) is a compactly supported \(1\)-form on an oriented surface
    \(S\) with boundary \(\partial S\), then \(\int_{\partial S}\alpha =
    \int_Sd\alpha.\)
\end{theorem}

Note we haven't defined manifolds with boundary but a precise definition with
nice pictures can be found on Wikipedia. Locally, the boundary plays the same
role as the curve in Green's theorem. Since the result is local, we can reduce
to this case.

Thus on an oriented surface \(S\), we have spaces of \(0,1,2\)-forms
\(\Omega^k(S)\) and exterior derivative operators \(d:\Omega^k(S)\to
\Omega^{k+1}(S)\). We also know how to integrate \(k\)-forms on
\(k\)-dimensional submanifolds (\(k\geq 1\)), and Stoke's theorem which
relates the two integrals. Lastly, we have the wedge product
\(\wedge:\Omega^1(S)\times \Omega^1(S)\to \Omega^2(S)\).

\subsection{The de Rham complex}
Let \(S\) be a smooth surface, and consider the sequence of maps
\[
    0\to \Omega^0(S)\xrightarrow{d} \Omega^1(S)\xrightarrow{d} \Omega^2(S)\to 0.
\]
We know \(d\circ d=0\) so this is a chain complex, and we can define the \textit{de Rham cohomology groups} as the cohomology spaces of this complex explicitly given by
\begin{align*}
    \HH^0(S) &= \ker(d:\Omega^0(S)\to \Omega^1(S)) \\
    \HH^1(S) &= \frac{\ker(d:\Omega^1(S)\to \Omega^2(S))}{\img(d:\Omega^0(S)\to
    \Omega^1(S))}\\
    \HH^2(S) &= \frac{\Omega^2(S)}{\img(d:\Omega^1(S)\to \Omega^2(S))}.
\end{align*}

We say a \(k\)-form \(\alpha\) is \textit{closed} if \(d\alpha=0\) and
\textit{exact} if \(\alpha=d\beta\) for some \((k-1)\)-form \(\beta\). Then every
exact form is closed and the de Rham cohomology measures `failure of closed
forms to be exact'.

In general, \(H^0(S)=\mathbb{R}\) for a connected surface \(S\). We will now
compute some examples.

\begin{lemma}[(de Rham cohomology of \(\mathbb{R}^2\))]
    We have \(\HH^0(\mathbb{R}^2)=\mathbb{R}\), and all other cohomology
    groups vanish.

    \begin{proof}
        Since \(\mathbb{R}^2\) is connected, we know
        \(\HH^0(\mathbb{R}^2)=\mathbb{R}\). The claim \(\HH^1(\mathbb{R}^2)=0\)
        follows from the criterion for a closed \(1\)-form to be exact. Thus it
        remains to show \(\HH^2(\mathbb{R}^2)=0\). Take a closed \(2\)-form
        \(\rho=R(x_1,x_2)dx_1\wedge dx_2\) and claim we can find functions
        \(\alpha_1,\alpha_2\) with \(R(x_1,x_2)=\frac{\partial
        \alpha_1}{\partial x_2}-\frac{\partial \alpha_2}{\partial x_1}\), which
        will show the required result. This can be done, for example, by setting
        \(\alpha_2=0\) and \(\alpha_1(x_1,x_2)=\int_0^{x_2} R(x_1,t)dt\).
    \end{proof}
\end{lemma}

\begin{proposition}
    Let \(S=\mathbb{S}^2\) be the sphere. Then \(\HH^1(S)=0\).
    \begin{proof}
        Suppose \(\alpha\) is a closed \(1\)-form on \(S\). Take the standard
        charts \(U,V\subseteq S\) given by removing the north and south
        poles from \(S\) respectively. Then \(U,V\cong \mathbb{R}^2\) so we can
        find functions \(f_U\), \(f_V\) on \(U,V\) respectively such that
        \(df_U=\alpha|_U\), \(df_V=\alpha|_V\). On \(U\cap V\), we see that
        \(f_U-f_V\) is locally constant hence is constant. Thus adding a
        constant function to \(f_V\) if necessary, we see that \(f_U\)
        and \(f_V\) agree on \(U\cap V\) and hence define a global
        function \(f\in \Omega^0(S)\) with \(df=\alpha\).
    \end{proof}
\end{proposition}

\begin{proposition}
    Let \(T=\mathbb{S}^1\times \mathbb{S}^1\) be the torus. Then \(\HH^1(T)\) is
    a two-dimensional vector space.

    \begin{proof}
        Consider the standard angular coordinates \(\vartheta, \varphi\in
        [0,2\pi)\). Let \(\gamma_1,\gamma_2\) be the standard circles
        corresponding to \(\vartheta=0\) and \(\varphi=0\) respectively, and
        consider the map
        \begin{align*}
            F:\Omega^1(T)\to \mathbb{R}^2, \qquad
                \alpha\mapsto \left(\int_{\gamma_1}\alpha,
                \int_{\gamma_2}\alpha\right).
        \end{align*}
        This is linear and vanishes on exact \(1\)-forms by the fundamental
        theorem of calculus, so we get a well-defined map \(F:\HH^1(T)\to
        \mathbb{R}^2\). We will show \(F\) is injective and surjective.

        If \(\alpha\) is a closed \(1\)-form such that
        \(\int_{\gamma_1}\alpha=\int_{\gamma_2}\alpha = 0\), then writing
        \(\alpha=Pd\vartheta + Qd\varphi\) locally shows that for any fixed
        \(\varphi\) we have
        \[
            \int_{\gamma_2}\alpha = \int_0^{2\pi} P(u,
            \varphi)du = 0
        \]
        and hence the function \(f(\vartheta, \varphi)=\int_0^\vartheta
        P(u,\varphi)du\) is well-defined and satisfies \(\frac{\partial
        f}{\partial \vartheta}=P(\vartheta,\varphi)\). Thus
        \(\tilde{\alpha}=\alpha-df\) is a closed \(1\)-form of the form
        \(Q\cdot d\varphi\). But \(d\tilde\alpha=0\) implies \(Q\) is
        constant and we see that \(Q=0\) since  \(\int_{\gamma_2}Qd\varphi=0\).
        Thus \(\alpha=df\) is exact.

        To show surjectivity, we consider the \(1\)-forms \(d\vartheta\),
        \(d\varphi\) which get mapped to \((1,0)\) and \((0,1)\)
        respectively.
    \end{proof}
\end{proposition}

\begin{remark}
    In general, a similar argument shows
    \(\HH^1(S)=\Hom(\Pi_1(S,p), \mathbb{R})\)
    where \(\Pi_1(S,p)\) is the fundamental group of \(S\) for some base point
    \(p\). In particular if \(\Sigma_g\) is the closed orientable surface of
    genus \(g\) then \(\HH^1(\Sigma_g)\cong \mathbb{R}^{2g}\).
\end{remark}

\subsection{Compactly supported cohomology}
We next calculate \(\HH^2(S)\). For this, we need to first modify the definition
of de Rham cohomology when \(S\) is non-compact, by setting
\(\Omega^i_C(S)\subseteq \Omega^i(S)\) to be the subspace of compactly supported
\(i\)-forms. The exterior derivative respects the compact support property,
hence we can define \(\HH^i_C(S)\) to be the \(i\)-th cohomology of the complex
\[
    0\to \Omega^0_C(S)\xrightarrow{d} \Omega^1_C(S)\xrightarrow{d}
    \Omega^2_C(S)\to 0.
\]
The main theorem is the following.

\begin{theorem}
    For any connected surface \(S\), we have \(\HH^2_C(S)=\mathbb{R}\).
    \begin{proof}
        We prove this in two steps-- first the special case \(S=\mathbb{R}^2\),
        and then use this to handle general surfaces.

        \underline{In the case of \(\mathbb{R}^2\)}, the proof is by explicit
        integration.  We will consider the map \(I:\Omega^2_C(\mathbb{R}^2)\to
        \mathbb{R}\) given by \(I(\rho)=\int_{\mathbb{R}^2} \rho\). This is
        well-defined by the compact support property. Surjectivity follows since
        we can explicitly find a smooth function \(\psi\) with \(I(\psi\cdot
        dx_1\wedge dx_2)=\int_{\mathbb{R}^2}\psi \cdot dx_1dx_2=1\).
        \marginpar{Lecture 4, 31/10/23}
        Thus it remains to show that the kernel of \(I\) is
        precisely the space of exact forms. One containment is obvious by
        Stokes' theorem, so suppose \(I(\rho)=0\) and write
        \(\rho=R(x,y)dx\wedge dy\) where (up to translation) \(R\) is supported
        in some box \([0,N]^2\). We construct \(\alpha\in
        \Omega^1_C(\mathbb{R}^2)\) with \(d\alpha=\rho\) by integrating \(R\)
        cleverly.

        First consider the `strong special case' when \(\int_{0}^N R(x,y)dx
        =0\) for all \(y\). Define \(g(x,y)=\int_{0}^x R(t,y)dt\). Then
        \(g(N,y)=0\) so \(g(x,y)dy\) is a compactly supported \(1\)-form and its
        derivative is \(R(x,y)dx\wedge dy\)  as required. In general, consider
        \(g(x,y)dy\). Then to reduce to the strong special case, we must manage
        the quantity \(H(y)=\int_0^N R(x,y)dx\) which doesn't vanish in general. Take
        \(b:\mathbb{R}\to \mathbb{R}\) a `bump function' with compact
        support such that \(\int_{0}^Nb(x)dx=1\). Since \(H:\mathbb{R}\to
        \mathbb{R}\) also has compact support, so does \(h(x,y)=b(x)H(y)\). Moreover,
        \[\int_{\mathbb{R}^2} h(x,y)dx\wedge dy = \int_{0}^N H(y)dy =
        \int_{\mathbb{R}^2}R(x,y)dx\wedge dy.\] Thus  \(\int_{0}^N
        (R(x,y)-h(x,y))dx=0\), i.e.\ the \(2\)-form
        \((R(x,y)-h(x,y))dx\wedge dy\) satisfies the hypotheses of the strong
        special case and is hence exact. So to complete the proof, we need to show \(h(x,y)dx\wedge dy =d\beta\) for some compactly supported
        \(1\)-form \(\beta\). But reversing \(x,y\) we know that
        \(\int_{0}^Nh(x,y)dy = b(x)\int_{0}^N H(y)dy = 0\), so this again
        is the strong special case.

        Thus we wrote \(\rho\) as a sum of two \(2\)-forms, both of which
        satisfy the hypotheses of the strong special case and are hence exact.

        \underline{For a general connected surface \(S\)}, again consider the
        map \(I:\Omega^2_C(S)\to \mathbb{R}\) given by \(\rho \mapsto
        \int_S\rho\).  Surjectivity and vanishing on exact forms follows
        similarly as before, so we must show that \(\ker I\) is precisely the
        space of exact forms. Suppose \(\rho\in \Omega^2_C(S)\) is such that
        \(\int_S \rho =0\). By the compact support property, there are finitely
        many contractible charts \(U_1,\ldots,U_n\) such that \(\rho\) is supported
        on a connected open set \(U_1\cup U_2\cup \ldots \cup U_n\). Since
        we know the result for \(\mathbb{R}^2\), we know it holds for each
        \(U_i\), so we show the result by induction on \(n\).

        Suppose \(n>1\). By connectedness of \(\bigcup_i U_i\) we can arrange
        that \(U_i\cap U_{i+1}\neq \emptyset\) for all \(i\), so that
        \(V_1=U_1\cup \ldots\cup U_{n-1}\) and \(V_2=U_n\) are connected opens.
        Then by induction hypothesis, any \(2\)-form \(\tilde\rho\) compactly
        supported on \(V_i\) satisfying \(\int_S\tilde\rho=0\) must be exact. So
        take a partition of unity \(f_1,f_2\) such that
        \(\text{Supp}(f_i)\subset V_i\) and \(f_1+f_2=1\) on
        \(\text{Supp}(\rho)\). Then \(\rho = \rho_1 +\rho_2\), where
        \(\rho_i=\rho\cdot f_i\) supported on \(V_i\) . From \(\int_S\rho=0\),
        we obtain a constant \(c=\int_S \rho_1=-\int_S\rho_2\). Choose a
        \(2\)-form \(\sigma\) supported on \(V_1\cap V_2\) such that
        \(\int_S\rho = c\). Then setting \(\tilde\rho_1=\rho_1-\sigma\) and
        \(\tilde\rho_2=\rho_2+\sigma\), we see that by induction hypothesis
        \(\tilde\rho_i=d\alpha_i\) for some \(\alpha_1,\alpha_2\). Then for the
        \(1\)-form \(\alpha=\alpha_1+\alpha_2\), we have \(\rho=d\alpha\) as
        required.
    \end{proof}
\end{theorem}

\section{Complex differential forms}
In the previous section, \(S\) was a smooth surface with no additional
structure. At a point \(p\in S\) the cotangent space \(T_p^\ast S\) is defined so that a real-valued function \(f:S\to \mathbb{R}\) induces a cotangent vector \([df]_p\) (seen as an \(\mathbb{R}\)-linear map \(T_pS\to \mathbb{R}\)). Likewise, we can consider the \textit{complexified cotangent space} \(T_P^\ast
S^\mathbb{C}=\Hom_\mathbb{R}(T_PS,\mathbb{C})\) defined so that a complex-valued
function \(f:S\to \mathbb{C}\) induces a cotangent vector \([df]_P:T_pS\to \mathbb{C}\). This makes sense for any surface \(S\).

We now specialise to Riemann surfaces, where there is a distinguished class of
holomorphic functions. What does this extra structure buy us?

\begin{definition}
    Let \(V\) be a real vector space. A \textit{complex structure} on \(V\) is
    an \(\mathbb{R}\)-linear map \(J:V\to V\) such that \(J^2=-\text{Id}\).
    Given such a vector space with complex structure, we say an
    \(\mathbb{R}\)-linear map \(A:V\to \mathbb{C}\) is \textit{complex linear}
    if \(A(Jv)=i\cdot A(v)\) for all \(v\in V\) and it is \textit{complex
    anti-linear} if
    \(A(Jv)=-i\cdot A(v)\) for all \(v\in V\).
\end{definition}

The cotangent spaces of a Riemann surface are naturally equipped with a complex
structure.

\begin{lemma}
    Given a Riemann surface \(S\) and a point \(p\in S\), there is a unique way
    to define a complex structure on \(T_pS\) such that the
    differential of any holomorphic function is complex linear.
    \begin{proof}[Sketch proof]
        This is a consequence of the fact that on open sets in \(\mathbb{C}\),
        the map sending a holomorphic function to its derivative is
        \(\mathbb{C}\)-linear. Choosing a local complex coordinate \(z=x+iy\)
        around \(p\), the complex structure \(J:T_pS\to T_pS\) then is given by
        multiplication by \(\sqrt{-1}\) i.e.\ a tangent vector determined by a
        path \(\gamma:\mathbb{R}\to S\) gets sent to the tangent vector
        \([i\cdot \gamma]\).
    \end{proof}
\end{lemma}

Let \(V\) be an \(\mathbb{R}\)-vector space with a complex structure.
Considering the \(\pm i\)-eigenspaces, we have the following useful piece of
linear algebra: any \(\mathbb{R}\)-linear map \(A:V\to \mathbb{C}\) can be
written as a sum of a complex linear map \(A'=\frac{1}{2}(A-i\cdot A\circ J)\)
and a complex anti-linear map \(A''=\frac{1}{2}(A+i\cdot A\circ J)\), thus
giving an orthogonal decomposition of \(\Hom_\mathbb{R}(V,\mathbb{C})\).

In the setting of Riemann surfaces, this gives a decomposition \(T_p^\ast
S^\mathbb{C}=T_p^\ast S'\oplus T_p^\ast S''\) of the complexified cotangent
space into a complex linear part (\textit{the holomorphic cotangent space}) and
an anti-linear part (\textit{the anti-holomorphic cotangent space})
respectively. The complex structure on \(T_pS\) was chosen such that for any
holomorphic function \(f\) we have \([df]_p\in T_p^\ast S'\) and
\([d\overline{f}]_p\in T_p^\ast S''\).


\subsection{Decomposing \(1\)-forms and exterior derivatives} The decomposition
of the cotangent space extends to the space of smooth complex \(1\)-forms naturally,
and we write \(\Omega^1(S,\mathbb{C})=\Omega^{1,0}(S)\oplus \Omega^{0,1}(S)\).
The exterior derivatives \(d:\Omega^0(S,\mathbb{C})\to
\Omega^1(S,\mathbb{C})\) and \(d: \Omega^1(S,\mathbb{C})\to
\Omega^2(S,\mathbb{C})\) decompose in a compatible fashion: there is a
commutative square
\[ \begin{tikzcd}[row sep=huge, column sep=huge]
        \Omega^{0,1}(S) \rar{-\partial (=-d)} & \Omega^2(S,\mathbb{C}) \\
        \Omega^0(S,\mathbb{C})\uar{\bar\partial} \rar{\partial} &
        \Omega^{1,0}(S)\uar{\bar\partial (=d)}
    \end{tikzcd}\]
where the \textit{Dolbeault operators} \(\partial, \bar\partial\) are given by
composing \(d\) with the natural projections \(\Omega^1(S,\mathbb{C})\to
\Omega^{0,1}(S), \Omega^{1,0}(S)\). In particular, we have \(d=\partial
+\bar\partial\) and \(\partial^2=\bar\partial^2=0\). The construction is explicit on a chart--
choosing a local complex coordinate \(z=x+iy\), the vectors \(dz=dx+idy\) and \(d\bar z= dx-idy\) give bases of \(T_p^\ast S'\) and \(T_p^\ast
S''\) respectively for \(p\) in the chart. Given a complex function \(f\) on this chart, we can write
\[
    df = \underbrace{\frac{1}{2}\left( \frac{\partial f}{\partial
       x}-i\frac{\partial f}{\partial y}\right)}_{\frac{\partial
       f}{\partial z}}dz +
       \underbrace{\frac{1}{2}\left( \frac{\partial f}{\partial
       x}+i\frac{\partial f}{\partial y}\right)}_{\frac{\partial
       f}{\partial \bar z}}d\bar z
\]
which tells us how to define the operators \(\partial\) and \(\bar\partial
\) on \(\Omega^0_{S,\mathbb{C}}\). Explicitly, we have \(\partial f =
\frac{\partial f}{\partial z}dz\) and \(\bar\partial = \frac{\partial f
}{\partial \bar z}d\bar z\). Likewise, computing the derivative
\(d:\Omega^1(S,\mathbb{C})\to \Omega^2(S,\mathbb{C})\) locally gives
\begin{align*}
    d(f\cdot dz +g\cdot d\bar z)&=\bar\partial(f\cdot dz)+\partial(g\cdot d\bar
    z)\\ &= \left(\frac{\partial g}{\partial z}-\frac{\partial f}{\partial \bar
    z}\right)dz\wedge d\bar z.
\end{align*}


\subsection{Holomorphic \(1\)-forms} A formal consequence of the definitions is that
\(\bar \partial f=0\) is locally the Cauchy--Riemann equation, thus the subspace
of holomorphic functions in \(\Omega^0(S,\mathbb{C})\) is precisely \(\ker
\bar\partial\).  In particular for a holomorphic function \(f\), we have
\(df=\partial f=f'\cdot dz\) locally where \(f'\) is the usual complex
derivative.

We extend this idea to define holomorphic \(1\)-forms, which will be holomorphic
complex linear forms.

\begin{definition}
    A \textit{holomorphic \(1\)-form} is a \((1,0)\)-form satisfying
    \(\bar\partial \beta\)=0.
\end{definition}

Locally if \(\beta=f\cdot dz\), then \(\beta\) is holomorphic if and only if \(f\) is holomorphic.

Note for a holomorphic \(1\)-form \(\beta\) we have \(d\beta=\partial \beta +
\bar\partial \beta =  0\), so holomorphic \(1\)-forms are always closed.
This gives an analogue to Cauchy's theorem: if \(Q\subset S\) is a surface
with boundary, then by Stokes' theorem we have \[\int_{\partial Q}\beta = 0.\]  Thus the integral of a holomorphic \(1\)-form vanishes on any closed
contour. The `meromorphic' analogue is as follows-- let \(\alpha\) be a
meromorphic \(1\)-form on \(S\), i.e.\ a \(1\)-form that is holomorphic on the
complement of a discrete subset \(D\subset S\) and is given locally as
\(\alpha=f\cdot dz\) for some meromorphic function \(f\) (one checks this notion
is well-defined independent of chart). Choosing \(D\) minimal, the points of
\(D\) give the poles of \(\alpha\). Let \(p\) be a pole, \(C\) a small loop
around \(p\).

\begin{definition}
    The residue of \(\alpha\) at \(p\) is
    \[\text{Res}_p(\alpha)=\frac{1}{2\pi i}\int_C \alpha\]
    which is the usual residue in local coordinates (i.e.\ the coefficient of
    \(z^{-1}dz\) in a Laurent expansion of \(\alpha\).)
\end{definition}

\begin{proposition}
    Let \(\alpha\) be a meromorphic \(1\)-form on a compact Riemann surface
    \(S\). Then the sum of the residues of \(\alpha\) over all poles is zero.

    \begin{proof}
        By compactness, the set of poles is finite say \(p_1,\ldots,p_n\). Around
        each \(p_j\), choose a small disc \(D_j\) with boundary given by a loop
        \(C_j\). We have
        \[\sum_{j=0}^n \text{Res}_{p_j}(\alpha)= \sum_{j=0}^n \left(
        \frac{1}{2\pi i}\int_{C_j}\alpha \right) =\frac{1}{2\pi
        i}\int_{\bigcup_j C_j} \alpha.\]
        Set \(\tilde S\) to be the complement of \(D_1\cup \ldots \cup D_n\). Then
        the boundary of \(\tilde{S}\) is \(-\bigcup_j C_j\) and \(\alpha\) is a
        holomorphic \(1\)-form on \(\tilde S\), so we have
        \(\int_{\bigcup_jC_j}\alpha = -\int_{\partial\tilde S}d\alpha = 0\).
    \end{proof}
\end{proposition}

\subsection{The Laplacian}
Since \(d^2=(\partial +\bar\partial)^2=0\), we have \(\partial \bar\partial =
-\bar\partial \partial\). This gives a canonical operator \(\Omega^0_S\to
\Omega^2_S\) on a Riemann surface.

\begin{definition}
    The \textit{Laplacian} is defined to be \(\Delta=2i\cdot
    \partial\bar\partial \). We say a function \(f\) is \textit{harmonic} if \(\Delta f=0\).
\end{definition}

Choosing a local complex coordinate \(z=x+iy\), we have \(\Delta f =
-\frac{1}{2}(\frac{\partial^2 f}{\partial x^2} + \frac{\partial^2 f}{\partial
y^2})dx\wedge dy\) so this is just the usual Laplacian (an operator that takes
smooth functions to smooth functions) multiplied by a chosen volume-form. The
presentation using differential operators removes the choice of this volume
form and makes the operator canonical.

If \(f\) is holomorphic, then its real and imaginary parts are harmonic.
Indeed, writing \(\text{Re}(f)=\frac{1}{2}(f+\bar f)\), we see that
\(\partial\bar\partial (f+\bar f)=\partial \bar \partial  f +
\overline{\partial \bar\partial f}=0\), likewise for the imaginary part.

\begin{exercise}
    If \(\varphi\) is real valued and defined on a neighbourhood \(U\subset
    S\) of \(p\), show that there is a \(U'\subset U\) such that
    \(\varphi|_U'\) is the real part of a harmonic function on \(U'\). Thus
    every real function has a corresponding imaginary `harmonic conjugate'.
\end{exercise}

\begin{exercise}[Maximum principle]
    Suppose \(\varphi\) is non-constant, real valued, harmonic on a
    connected open set \(U\subset S\). Then show that for any \(x\in U\) there
    is a point \(x'\in U\) with \(\varphi(x')>\varphi(x)\).
\end{exercise}

\begin{theorem}[(Inverting the Laplacian)]
    \label{thm:laplace-inversion}
    Suppose \(\rho\) is a \(2\)-form on a compact surface \(S\). Then
    \(\int_S\rho=0\)if and only if there is a smooth function \(f\) with
    \(\Delta f=\rho\), and if so \(f\)is unique up to the addition of a
    constant.
\end{theorem}

This is the main technical result of the course. We will use it to prove the
uniformisation theorem and the Riemann-Roch theorem, which in turn can be used
to show that all compact Riemann surfaces are algebraic.

\begin{remark}
    The result on inverting the Laplacian really applies to all elliptic linear
    partial differential operators on compact manifolds. In this context, the
    result is often called the `Fredholm alternative'.
\end{remark}

\section{Topological applications}

We will now study the applications of the theory developed-- in particular the
Euler characteristic, the Riemann--Hurwitz theorem, and Dolbeault cohomology.
\marginpar{Lecture 5, 07/11/23} We will later also look at analytic
applications, solving Laplace's equation on a Riemann surface.

\subsection{The Euler characteristic} This is a purely topological notion. We
will use some techniques from differential topology without proving them, and
this will result in the lecture being a little sketchier. The goal is to
understand the topological genus of a Riemann surface through differential
forms.

Let \(S\) be a smooth surface. We can triangulate it to define the Euler
characteristic as
\[
    \chi(S)=\#\{\text{vertices}\}-\#\{\text{edges}\}+\#\{\text{faces}\}.
\]
The important fact is that this is independent of the triangulation chosen;
one proves this by showing that the Euler characteristic is invariant under
refinement (subdivision) of the triangulation and that any two triangulations
admit a common refinement.

By explicit triangulation, one can calculate that the Euler
characteristic of \(\Sigma_g\) (the compact surface of genus \(g\)) is
\(2-2g\). In fact this is precisely how we define genus, i.e.\ for a surface
\(S\) we have
\[
    g(S)=1-\frac{1}{2}\chi(S).
\]

But the genus appears in other forms, for example in the first de Rham
cohomology as we have seen previously. We will see how the two notions are
related. Suppose \(S\) is oriented, and let
\(\alpha\) be a \(1\)-form such that the set \(\Delta\subset S\) where
\(\alpha\) vanishes is discrete. For any \(p\in \Delta\), we can choose local
coordinates \(x_1,x_2\) such that \(p=(0,0)\) and
\(\alpha=\alpha_1dx_1+\alpha_2dx_2\). Thus for small \(r>0\), the only zero of the function
\[
    (\alpha_1,\alpha_2): \overline{B(0,r)}\to \mathbb{R}^2
\]
is at the origin. In particular, when restricted to the boundary,
\((\alpha_1,\alpha_2)|_{\partial B(0,r)}\) gives a curve in
\(\mathbb{R}^2\setminus 0\) which has a well-defined winding number (the
signed number of times the curve wraps around the origin anticlockwise). One
checks that this is independent of local coordinates (diffeomorphism invariance)
and independent of choice of \(r\) (homotopy invariance) since it is a
continuous integer-valued function. Thus we can make sense of the
\textit{winding number of a \(1\)-form around a point}, also called the
\textit{index of the \(1\)-form at the point}.  Write
\(\nu_p(\alpha)\) for this number, noting that it can be computed in various ways for example by integrating \(d\theta\) around \((\alpha_1,\alpha_2)\partial
B(0,r)\)  in local polar coordinates.

Note that this notion makes sense even for \(p\notin \Delta\), but if
\(\alpha(p)\neq 0\) then we can choose \(r\) sufficiently small so that
\((\alpha_1,\alpha_2):\overline{B(0,r)}\to \mathbb{R}^2\) does not hit the
origin. Thus the winding number of the boundary curve around the origin is
necessarily zero.

\begin{proposition}[(Poincar\'e--Hopf theorem for smooth \(1\)-forms)]
    For \(S, \alpha, \Delta\) as above, the Euler characteristic of \(S\) can be
    computed as \[\chi(S)=\sum_{p\in \Delta} \nu_p(\alpha).\]
    \begin{proof}[Very sketchy proof]
        One could (dually, equivalently) use vector fields rather than
        \(1\)-forms by fixing a volume form. The result is then the usual
        Poincar\'e--Hopf theorem\footnote{One of the consequences of the
        Poincar\'e--Hopf theorem is the hairy ball theorem.}.

        One first uses the Gauss map to show that the number \(\sum_{p\in
        \Delta}\nu_p(\alpha)\) is indepentent of \(\alpha\), see [Milnor,
        \textit{"Topology from the differentiable viewpoint"}, Chapter 6]. This
        lets us choose a specific \(1\)-form whose winding number we can relate
        to the triangulation. In particular, given a triangle we choose a
        function \(f\) definedon a neighbourhood which has minima on the three
        vertices, maxima in the interior of the face, and saddle points on the
        edges. The zeroes of \(df\) are the critical points of \(f\), precisely
        the seven points described above. Then the winding number of \(df\)
        around a minimum or a maximum is \(+1\) by an explicit calculation
        (using Taylor series, say), and at a saddle point it is \(-1\). One
        glues these functions to get a globally defined smooth function, and use
        its differential to get the required computation.
    \end{proof}
\end{proposition}

On a compact Riemann surface \(S\), we like holomorphic and meromorphic
\(1\)-forms.  To a holomorphic \(1\)-form \(\alpha\), we associate a real
\(1\)-form \(A=\frac{1}{2}(\alpha +\bar\alpha)\). Locally if \(\alpha=f(z)dz\),
then \(A= \frac{1}{2}(f(z)dz +\bar f(z)d\bar z)\) hence \(A\) vanishes if and
only if \(\alpha\) does. In particular \(A\) vanishes on a discrete set
\(\Delta\) so can be used to compute the Euler characteristic of \(S\).

We will show that the index of \(A\) at \(p\in S\) computes the vanishing
multiplicity of \(\alpha\) at \(p\) (defined by writing \(\alpha=f(z)dz\)
locally), written \(m_p(\alpha)\). This will give
the following version of the Poincar\'e--Hopf formula.

\begin{proposition}[(Poincar\'e--Hopf theorem for holomorphic \(1\)-forms)]
    For \(S\) a compact Riemann surface with a holomorphic \(1\)-form
    \(\alpha\) vanishing on a discrete set \(\Delta\), we have
    \[
        \sum_{p\in \Delta}m_p(\alpha) = -\chi(S).
    \]
    \begin{proof}
        Say \(f=f_1dz+if_2dz\) for smooth real functions \(f_1,f_2\), and local
        coordinate \(z=x+iy\). Then we see that \(A=f_1dx-f_2dy\), i.e.\ the
        winding number of \(A\) around \(p\in \Delta\) is the negative of the
        winding number of \((f_1,f_2)=f:\partial B(0,r)\to \mathbb{R}^2\). Now
        by the argument principle, this number \(-\nu_p(A)\)
         is equal to the vanishing multiplicity of \(f(z)\) at
        \(p\).  Thus the total number of zeros of \(\alpha\) with multiplicity
        is \(2g-2=-\chi(S)\).
    \end{proof}
\end{proposition}

Note this quantity has to be positive. Thus if \(g=0\) there are no holomorphic
\(1\)-forms at all. If \(g=1\), any holomorphic \(1\)-form would have to be
nowhere vanishing (and indeed such \(1\)-forms exist).

Holomorphic \(1\)-forms are too restrictive. Instead we consider meromorphic
\(1\)-forms. Fix a volume-form \(\omega\) on \(S\), locally given by
\(ig(z)dz\wedge d\bar z\) for a positive function \(g\). This gives a hermitian
metric on \(T^\ast S\): for a cotangent vector \(v\in T_p^\ast S\), we
define \(|v|^2\) so that \(v\wedge \bar v = |v|^2 \omega(p)\).  Note \(|v|^2\geq
0\), and it vanishes if and only if \(v=0\).

Let \(\alpha\) be a meromorphic \(1\)-form on \(S\). Choose a real-valued
function \(\rho\) on \(\mathbb{R}\) with \(\rho(t)=1\) for small \(t\) and
\(\rho(t)=t^{-1}\) for large \(t\). Define a new differential form
\(\tilde\alpha\) satisfying \(\tilde\alpha = \rho(|\alpha|^2)\alpha\) away from
the poles of \(\alpha\), and \(\tilde\alpha = 0\) at poles of \(\alpha\).
Then locally near a pole of \(\alpha\), we have
\[
    \tilde\alpha = \frac{1}{|f(z)|^2}f(z)R\cdot dz = \bar f(z)^{-1} R\cdot dz
\]
where \(R\) is a smooth positive function. In particular this is a smooth
\(1\)-form that vanishes at the poles. Thus the zero set of \(\tilde\alpha\) is
the set of zeros and poles of \(\alpha\). Moreover, at a pole of order \(d\) of
\(\alpha\), \(\tilde\alpha\) has index \(-d\) since \(\bar f\) is
anti-holomorphic. This gives us the following.

\begin{corollary}
    If \(\alpha\) is a nontrivial meromorphic \(1\)-form on a compact Riemann
    surface of genus \(g\), then the number of zeros of \(\alpha\) minus the
    number of poles of \(\alpha\) (counted with multiplicity) is \(2g-2\).
\end{corollary}

\subsection{Riemann--Hurwitz formula} The results of the previous section allow
for a quick proof of this classical result. Let \(f:X\to Y\) be a nonconstant
holomorphic map between compact Riemann surfaces. We defined a multiplicity
\(k_x\) associated to each \(x\in X\) such that locally the map is \(z\mapsto
z^{k_x}\), with \(k_x=1\) at all points except
for a finite set of ramification points.

\begin{definition}
    The \textit{ramification index} of \(f\) is defined to be
    \(
        R_f= \sum_{x\in X} (k_x-1),
    \)
    noting that the sum is finite.
\end{definition}

We also defined the degree \(d\) of \(f\), the number of points in any preimage
counted with multiplicity. These quantities are related as follows.

\begin{theorem}[(Riemann--Hurwitz)]
    The genera \(g_X\) of \(X\) and \(g_Y\) of \(Y\) are related as
    \[
        2-2g_Y = d(2-2g_X)-R_f.
    \]
    \begin{proof}
        We will prove this assuming that there is a meromorphic \(1\)-form on
        \(\beta\) on \(Y\). This is a harmless assumption, we will later use the
        Riemann--Roch theorem to prove that this always holds for every Riemann
        surface.

        Given this, \(f^\ast \beta\) is a meromorphic \(1\)-form on \(X\). If
        \(x\in X\) is a point around which \(f\) looks like \(z\mapsto z^k\),
        then writing \(\beta=g(w)dw\) locally around \(f(x)\) gives \(f^\ast
        \beta = kz^{k-1}g(z^k)dz\). Thus if \(\beta\) has a zero of order
        \(\ell\) at y=\(f(x)\), then \(f^\ast \beta\) has a zero of order
        \(k\ell + k -1\) at \(x\). So the contribution to the count of zeros and
        poles of \(f^\ast\beta\) coming from \(y\) is
        \[
            \sum_{x\in f^{-1}(y)}(k_x\ell + k_x - 1) = d\cdot \ell + \sum_{x\in
            f^{-1}y}(k_x-1).
        \]
        The result follows from the formula of Euler characteristic in terms of
        meromorphic \(1\)-forms.
    \end{proof}
\end{theorem}

\begin{exercise}
    Another application is the `degree--genus formula', which states that given
    a surface \(S\subset \mathbb{P}^2\) defined using a homogeneous degree \(d\)
    polynomial, we have \(g(S)=\frac{1}{2}(d-1)(d-2)\). As a corollary, one
    notes that not all Riemann surfaces can be embedded in \(\mathbb{P}^2\).
\end{exercise}

\subsection{Dolbeault cohomology} We continue to relate complex differential
forms to geometry. Recall that we used the exterior derivative
\(d:\Omega^q(S)\to \Omega^{q+1}(S)\) to define the de Rham cohomology.We now use
the operator \(\bar\partial:\Omega^{p,q}(S)\to \Omega^{p,q+1}(S)\) to define
\textit{Dolbeault cohomology}, where \(\Omega^{p,q}(S)\subset \Omega^{p+q}(S)\)
is the space of \((p,q)\)-forms on \(S\). For example, \(\Omega^{0,0}(S)\) is
the space of complex functions on \(S\), while \(\Omega^{0,1}(S)\) is the space
of \(1\)-forms locally given by \(fd\bar z\) for a complex function \(f\).

\begin{definition}
    The Dolbeault cohomology groups of \(S\) are
    \begin{align*}
        \HH^{0,0}(S)&=\ker(\bar\partial : \Omega^{0,0}(S)\to \Omega^{0,1}(S))\\
        &=\{\text{global holomorphic functions}\}\\
        \HH^{1,0}(S)&=\ker(\bar\partial : \Omega^{1,0}(S)\to \Omega^{1,1}(S))\\
        &=\{\text{global holomorphic \(1\)-forms}\}\\
        \HH^{0,1}(S)&=\frac{\Omega^{0,1}(S)}{\img(\bar\partial:\Omega^{0,0}(S)\to
        \Omega^{0,1}(S))}\\
        \HH^{1,1}(S)&=\frac{\Omega^{1,1}(S)}{\img(\bar\partial:\Omega^{1,0}(S)\to
        \Omega^{1,1}(S))}.
    \end{align*}
\end{definition}

\begin{exercise}
    For \(S= \mathbb{C}\), note that \(\HH^1(S)=0\) but
    \(\HH^{1,0}(S)\) is infinite dimensional. Now consider a \((0,1)\)-form on
    \(\mathbb{C}\) of the form \(fd\bar z\). Show that there is a complex
    function \(g\) on \(\mathbb{C}\) with \(\bar\partial g = fd\bar z\). So
    \(\HH^{0,1}(\mathbb{C})=0\). This is called the \(\bar\partial\)--Poincar\'e
    lemma.
\end{exercise}

Note that while we have good geometric interpretations of \(\HH^{0,0}\) and
\(\HH^{1,0}\), interpreting \(\HH^{0,1}\) takes a bit more work. Suppose we try
to find a meromorphic function on a Riemann surface \(S\) with a simple pole at
\(p\in S\) and no other poles. This would, for instance, show that \(S\cong
\mathbb{P}^1\). Finding such a function locally is trivial --- just pick a local
coordinate \(z\) in a neighbourhood \(U\ni p\) and choose the function
\(\frac{1}{z}\). Let \(\beta\) be a smooth cut-off function near \(p\)  (i.e.\ a
real function supported on \(U\) and equal to \(1\) near \(p\)). So \(\beta(\frac{1}{z})\) induces a smooth function on \(S\setminus \{p\}\). We
want a smooth function \(g\) on \(S\) such that \(g+\beta(\frac{1}{z})\) is
holomorphic on \(S\setminus \{p\}\). Now \(A=\bar\partial \beta(\frac{1}{z})\)
has compact support in \(U\setminus \{p\}\) (since \(\beta=1\) near \(p\)),
and we can extend \(A\) to a \((0,1)\)-form on \(S\) by extending by zero
over \(p\). So we want to solve \(\bar \partial g = -A\) for \(A\in
\Omega^{0,1}(S)\), and a solution exists if and only if the class \([A]\in
\HH^{0,1}(S)\) equals zero.

\begin{corollary}
    Suppose \(S\) is a compact Riemann surface with \(\HH^{0,1}(S)=0\). Then
    \(S\cong \mathbb{P}^1\).
\end{corollary}

\begin{theorem}
    Let \(S\) be a compact Riemann surface. We have the following operations and
    compatibilities on Dolbeault cohomology.
    \begin{enumerate}
        \item The map \(\sigma:\HH^{1,0}(S)\to \HH^{0,1}(S)\), taking a global
        holomorphic one-form \(\alpha\) to its conjugate \(\bar\alpha\), is a
        complex anti-linear isomorphism.
        \item The bilinear map \(B:\HH^{1,0}(S)\times\HH^{0,1}(S)\to
            \mathbb{C}\) given by \(B(\alpha, [\theta])=\int_S \alpha\wedge
            \theta\) is a perfect pairing. In particular there is an isomorphism
            \(\HH^{0,1}(S)\cong (\HH^{1,0}(S))^\ast\).
        \item The map \(\HH^{1,0}(S)\oplus \HH^{0,1}(S)\to \HH^1(S)\) defined by
        \((\alpha,\theta)\mapsto
        \iota(\alpha)+\overline{\iota(\sigma^{-1}\theta)}\) is an isomorphism,
        where \(\iota:\HH^{1,0}(S)\to \HH^1(S)\) sends a holomorphic \(1\)-form
        to its cohomology class.
        \item The map \(\gamma:\HH^{1,1}(S)\to \HH^2(S)\) defined by the natural
        inclusion \(\img(\bar\partial:\Omega^{0,1}\to \Omega^2)\subset
        \img(d:\Omega^1\to \Omega^2)\) is an isomorphism.
    \end{enumerate}
    \begin{proof}
        These are all consequences of inverting the Laplacian. We first show
        \(\sigma\) is surjective. Given \([\theta]\in \HH^{0,1}(S)\) we want a
        \(\theta'=\theta+\bar\partial f\) with \(\partial \theta'=0\), since
        then \(\alpha=\overline{\theta'}\) is a holomorphic \(1\)-form and
        \([\theta]=\sigma(\alpha)\). So we want to solve \(\partial
        \bar\partial f = -\partial \theta\). But since \(\int_S\partial
        \theta =0\), we can invert the Laplacian \(\partial
        \bar\partial=\frac{1}{2}i\Delta\) to produce \(f\).

        The composition of \(\sigma\) with \(B\) is, up to a factor, the
        Hermitian pairing \(\langle \alpha, \beta \rangle=\int_S
        \alpha\wedge \bar\beta\) which is positive definite. So \(B\) is a
        perfect pairing and \(\sigma\) is injective.

        The remaining two can be proved similarly, and are left as exercises.
    \end{proof}
\end{theorem}

We can now relate the topological genus to complex geometry.

\begin{corollary}
    For a compact Riemann surface \(S\) of genus \(g\), we have
    \(\textnormal{dim}\,\HH^{1,0}(S)=\textnormal{dim}\,\HH^{0,1}(S)=g\).
\end{corollary}

\section{Inverting the Laplacian}
\marginpar{Lecture 6, 14/11/23} Having seen the utility of
\cref{thm:laplace-inversion}, we will now repay the technical debt and look at
its proof. This will invoke key results from functional analysis which we will
recall as necessary.

\subsection{The Dirichlet norm} Let \(S\) be a Riemann surface,
\(\alpha\in \Omega^{1,0}(S)\) given locally by \(fdz\) for \(f\) a smooth
\(\mathbb{C}\)-valued function. Then \(i\alpha \wedge \bar\alpha\) is a
\((1,1)\)-form on \(S\) given locally as \(i\cdot |f|^2
dz\wedge d \bar z = |f|^2dx\wedge dy\). In particular this is a non-negative
real \(2\)-form hence we have \[||\alpha||^2 \coloneq \int_S i\alpha \wedge
\bar\alpha \geq 0.\] When \(\alpha\) has compact support, this quantity is
finite and thus defines a norm on the space of compactly supported
\((1,0)\)-forms \(\Omega^{1,0}_C(S)\). This norm is induced by the hermitian
inner product \(\langle \alpha,\beta \rangle = \int_S i \alpha \wedge \bar
\beta.\)

Fixing a volume-form (i.e.\ a nowhere vanishing real \(2\)-form) \(\omega\), we
can define a \textit{pointwise norm} \(|\alpha|^2\in \Omega^0(S)\),
characterised by the property \(i\alpha \wedge \bar \alpha = |\alpha|^2
\omega\). Then we have \(||\alpha||^2=\int_S |\alpha|^2 \omega,\) i.e.\ the norm
on \(\Omega^{1,0}_C(S,\mathbb{C})\) comes from the \(L^2\)-norm on
\(\Omega^0_C(S,\mathbb{R})\).  The norm on \((1,0)\)-forms is on the other hand
independent of the choice of volume-form.

Recall there is an isomorphism \(\Omega^{1,0}(S)\cong \Omega^{0,1}(S)\) given by
complex conjugation, and there is an orthogonal decomposition
\(\Omega^{1}(S,\mathbb{C})=\Omega^{1,0}(S)\oplus \Omega^{0,1}(S)\). Writing
\(\alpha^{1,0}\) and \(\alpha^{0,1}\) for the two components of a complex
\(1\)-form \(\alpha\) with respect to this decomposition, we can naturally
extend the norm above to \(\Omega^1_C(S,\mathbb{C})\) as \(||\alpha||^2 =
||\alpha^{1,0}||^2+||\alpha^{0,1}||^2\). Here \(||\alpha^{0,1}||\) is by
definition \(||\overline{\alpha^{0,1}}||\).

The space of compactly supported real \(1\)-forms is naturally the subspace
\(\{\alpha\in
\Omega^1_C(S,\mathbb{C})\;|\;\alpha^{1,0}=\overline{\alpha^{0,1}}\}\). Thus the
norm restricts, and \(\Omega^1_C(S,\mathbb{R})\) has an inner product and a norm given as
\[\langle A,B \rangle = 2i \int_S A^{1,0}\wedge B^{0,1}, \qquad ||A||^2 =
2||A^{1,0}||^2.\]
We can then use the exterior derivative to induce a norm on the space of
smooth real-valued functions.

\begin{definition}
    The \textit{Dirichlet inner product} on \(\Omega^0_C(S,\mathbb{R})\) is
    given by \(\langle f,g \rangle_D \coloneq \langle df,dg\rangle\). This
    induces the \textit{Dirichlet norm}, given by \(||f||_D = ||df||\). These
    definitions can also be made for the complex vector space
    \(\Omega^0(S,\mathbb{C})\) analogously.
\end{definition}

Note that the Dirichlet inner product (norm) is \underline{not} an inner product
(resp.\ norm) since, for instance, any non-zero constant function \(f\)
satisfies \(||f||_D=0\). Conversely if \(||f||_D=||df||\) vanishes, then
\(df=0\) since \(||\cdot ||\) is a norm on \(\Omega^1_C(S,\mathbb{C})\) and
hence \(f\) is constant. Thus it makes sense to consider the space
\(\Omega^0_C(S,\mathbb{R})/\mathbb{R}\) where we quotient out the subspace of
constant functions.

\begin{proposition}
    The Dirichlet inner product (norm) is an inner product (resp.\ norm) on
    \(\Omega^0_C(S,\mathbb{R})/\mathbb{R}\) and turns it into a pre-Hilbert
    space (i.e.\ a vector space equipped with an inner product).
\end{proposition}

We remark that the definition of the Dirichlet inner product makes sense even when just one of the arguments has compact support.

\begin{lemma}
    If \(f,g\) are smooth functions on \(S\) and \(f\) has compact support,
    then \(\langle f,g \rangle_D = \int_S f\cdot \Delta g = \int_S
    \Delta f \cdot g.\)

    \begin{proof}
        A partition of unity argument reduces this to the local case. We then
        have \(\langle f,g\rangle_D = 2i \int_S \partial f\wedge \bar\partial g
        \), so writing \(\partial f \wedge \bar\partial g = \partial
        (f\bar\partial g)-f\cdot\partial\bar\partial g \) and using Stokes'
        theorem gives the required result. Note this is just integration by
        parts.
    \end{proof}
\end{lemma}

We can now prove one direction of \cref{thm:laplace-inversion}.

\begin{corollary}
    Let \(S\) be a compact Riemann surface with a real \(2\)-form \(\rho\).
    Then up to adding a constant there is at most one \(f\in
    \Omega^0(S,\mathbb{C})\) satisfying \(\Delta f = \rho\), and such an \(f\)
    exists only if \(\int_S \rho = 0\).

    \begin{proof}
        To show uniqueness of the solution, we show \(\Delta f = 0\) (if and)
        only if \(f\) is a constant. Indeed we have \(\int_S f\Delta
        f = ||f||_D^2\), which vanishes if and only if \(df=0\) (i.e.\ \(f\) is
        constant) by properties of norms.

        If \(\rho = \Delta f\) for some function \(f\), then \(\int_S\rho =
        2i \int_S \partial \bar\partial f = 2i \int_S d(\partial f)\) which
        vanishes by Stokes' theorem.
    \end{proof}
\end{corollary}

\subsection{The heart of the proof} Our goal is to show that
a real \(2\)-form \(\rho\) on a compact Riemann surface \(S\) that satisfies
\(\int_S\rho = 0\) \ is in the image of \(\Delta\). Suppose \(\int_S\rho=0\),
and consider the pre-Hilbert space \(\Omega^0(S,\mathbb{R})/\mathbb{R}\).

\begin{exercise}
    For \(f\in \Omega^0(S,\mathbb{R})\), we have \(\Delta f = \rho \) if and
    only if for all \(g\in \Omega^0(S,\mathbb{R})\), \(\int_S g(\rho - \Delta
    f) = 0\).
\end{exercise}

Now we know for all \(f,g\) in this space, we have
\(\int_S\psi(\rho-\Delta\varphi) = \int_S\psi\rho - \langle
\varphi,\psi\rangle_D\). In particular, \(\Delta f = \rho\) if
and only if \(\int_S g\rho = \langle f,g\rangle_D\) for all \(g\). This reduces
the problem to studying the linear operator
\[
    \hat\rho :\Omega^0(S,\mathbb{R})/\mathbb{R} \to \mathbb{R}, \qquad g\mapsto
    \int_S g\rho
\]
which is well-defined since \(\int_S \rho = 0\). \underline{Claim} this is a
bounded linear operator on \(\Omega^0(S,\mathbb{R})/\mathbb{R}\) and hence
extends to the metric completion \(H =
(\Omega^0(S,\mathbb{R})/\mathbb{R})^{\wedge}\) which is the Hilbert space whose
points are given by equivalence classes of Cauchy sequences in
\(\Omega^0(S,\mathbb{R})/\mathbb{R}\).

Then the Riesz representation theorem states that every bounded linear operator
on \(H\) is of the form \(\langle \cdot, f\rangle_D\) for some \(f\in H\), thus
giving a weak solution to the problem-- we have an \(f\in H\) such that
\(\hat\rho(g)=\langle f, g\rangle_D\) for all \(g\). Finally, we
\underline{claim} that \(f\) is in fact in
\(\Omega^0(S,\mathbb{R})/\mathbb{R}\), i.e.\ \(f\) is smooth. This finishes the
proof of \cref{thm:laplace-inversion}.

\subsection{The boundedness claim} We first prove an analytic result about
smooth functions on \(\mathbb{R}^2\). Let \(U\subset \mathbb{R}^2\) be a bounded
convex open set of diameter \(d\) and area \(A\) with respect to the standard
Lebesgue measure \(d\mu\). For a smooth function \(\varphi\) defined on an open
set containing \(\bar U\), write \(\bar\varphi = \frac{1}{A} \int_U \varphi
d\mu\) for its average. We show the deviation of \(\varphi\) from its average is
bounded.

\begin{lemma}[(Poincar\'e inequality)]
    For \(U,d,A,\varphi\) as above and \(x\in U\), we have
    \[
        |\varphi(x)-\bar\varphi| \leq \frac{d^2}{2A}\int_{y\in
        U}\frac{|\nabla \varphi(y)|}{|x-y|} d\mu.
    \]

    \begin{proof}
        Up to translation and adding a constant, we can assume \(x=\mathbf{0}\)
        is the origin and \(\varphi(\mathbf{0})=0\) so we have to bound
        \(|\bar\varphi|\). Since \(U\) is convex, we
        can use polar coordinates to write \(\bar\varphi =
        \frac{1}{A}\int_0^{2\pi}\int_0^{R(\theta)}\varphi(r,\theta)rdrd\theta\).
        Moreover, we can write \(\varphi(r,\theta)=\int_0^r \frac{\partial
        \varphi}{\partial s}(s,\theta)ds\) since \(\varphi(\mathbf{0})=0\). Thus
        we have
        \begin{align*}
            \bar\varphi
            &= \frac{1}{A}\int_{\theta=0}^{2\pi} \int_{r=0}^{R(\theta)}
            \int_{s=0}^r \frac{\partial \varphi}{\partial s}(s,\theta)
            rdsdrd\theta\\
            &= \frac{1}{A} \int_{\theta=0}^{2\pi} \int_{s=0}^{R(\theta)}
            \left( \int_{r=s}^{R(\theta)} rdr\right) \frac{\partial
            \varphi}{\partial s}(s,\theta) dsd\theta
        \end{align*}
        and \(\int_{r=s}^{R(\theta)}rdr =\frac{1}{2}(R(\theta)^2-s^2) \leq
        \frac{1}{2}R(\theta)^2 \leq \frac{1}{2}d^2\). Additionally, we know
        \(|\frac{\partial \psi}{\partial s}|\leq |\nabla \varphi|\) so putting
        these together, we have the required inequality
        \begin{align*}
            |\bar\varphi|
            &\leq \frac{d^2}{2A} \int_{\theta=0}^{2\pi} \int_{s=0}^{R(\theta)}
             |\nabla\varphi(s,\theta)| \cdot \frac{1}{s} sdsd\theta \\
            &\leq \frac{d^2}{2A} \int_{y\in U} \frac{|\nabla
            \varphi(y)|}{|y|}d\mu.
        \end{align*}
    \end{proof}
\end{lemma}

\begin{corollary}
    \label{cor:poincare-ineq}
    Under the hypotheses above, we have \[\int_U
    |\varphi-\bar\varphi|^2d\mu \leq \frac{d^2\pi}{2A}\int_U|\nabla
    \varphi|^2 d\mu.\]
    \begin{proof}
        This is essentially a consequence of Yang's inequality on convolutions--
        for functions \(g,h\) on \(\mathbb{R}^2\) such that \(g\) is
        integrable and \(h\) is square-integrable, we have
        \[
            ||g\ast h||_{L^2} \leq ||g||_{L^1} \cdot ||h||_{L^2}
        \]
        where \(g\ast h(x)=\int_{y\in \mathbb{R}^2} g(y)h(x-y)d\mu\) is the
        \textit{convolution} of \(g\) and \(h\), and \(||\cdot ||_{L^1},
        ||\cdot ||_{L^2}\) denote the \(L^1\) and \(L^2\) norms respectively.
        Convolution \(\ast\) is a commutative and associative operation
        whenever defined.

        To prove the corollary we choose \(g(x)=\frac{d^2}{2A}\frac{1}{|x|}\) if
        \(0<|x|<d\) and \(0\) otherwise, noting this is
        discontinuous but nonetheless integrable with
        \(||g||_{L^1}=\frac{d^3\pi}{A}.\) Likewise, choose
        \(h(x)=|\nabla \varphi(x)|^2\) for \(x\in U\) and \(0\)
        otherwise. Then \(g\ast h\) is a positive function on
        \(\mathbb{R}^2\) and the Poincar\'e inequality says
        \(|\varphi(x)-\bar\varphi|\leq |g\ast h(x)|\) for all
        \(x\in U\). In particular
        \(\int_U|\varphi-\bar\varphi|^2d\mu\leq ||g\ast h||_{L^2}^2\), so Yang's inequality gives the result.
    \end{proof}
\end{corollary}

We can now show \(\hat\rho:\Omega^0(S,\mathbb{R})/\mathbb{R}\to \mathbb{R}\)
given by \(g\mapsto \int_S g\rho\) is a bounded operator. We first reduce to the
case when \(\rho\) is supported in a single coordinate chart-- since
integration over \(S\) defines an isomorphism \(H^2(S)\cong \mathbb{R}\), we
know \( [\rho]= \int_S\rho =0\) i.e.\ \(\rho=d\theta\) for some \(1\)-form.
By compactness of \(S\), we can use a partition of unity argument to write
\(\theta = \theta_1+\ldots+\theta_n\) where each \(\theta_i\) is a \(1\)-form
supported on a bounded convex coordinate chart. Then each \(\rho_i=d\theta_i\)
is a \(2\)-form supported on such a chart, satisfies \(\int_S\rho_i = 0\), and
we have \(\rho=\rho_1+\ldots+\rho_n\). Thus \(\hat\rho =
\hat\rho_1+\ldots+\hat\rho_n\), and it suffices to show each operator
\(\hat\rho_i\) is bounded.

Assume \(\rho\) is a \(2\)-form supported on a bounded convex coordinate chart
\(U\subset \mathbb{R}^2\). Using the standard nowhere-vanishing volume form
\(dx\wedge dy\), we can identify \(2\)-forms with functions so that
\(\rho=hdx\wedge dy\) for a smooth function \(h\) supported on \(U\). A function \(\varphi \in \Omega^0(S)\) induces a function
on \(U\) by restriction, and we have
\[
    |\hat\rho(\varphi)|
    = \bigg\vert \int_U h\varphi d\mu \bigg\vert
    = \bigg\vert \int_U h(\varphi -\bar \varphi) d\mu \bigg\vert
    \leq ||h||_{L^2(U)} \cdot ||\varphi-\bar\varphi||_{L^2(U)}
\]
where we recall \(\int_U h\bar\varphi d\mu = \bar\varphi \int_U hd\mu = 0\) and
use the Cauchy--Schwartz inequality. Thus the corollary gives us
\[|\hat\rho(\varphi)|\leq \frac{d^2\pi}{2 A} \cdot ||h||_{L^2(U)}\cdot ||\nabla
\varphi||_{L^2(U)}.\] Lastly, note that
\(||\nabla \varphi||_{L^2(U)}\leq ||d\varphi||_{L^2(X)}=||\varphi||_D\) and
write \(C=\frac{d^2\pi}{2 A}\cdot ||h||_{L^2(U)}\)  so that we have the required
inequality
\[
    |\hat\rho(\varphi)|\leq C\cdot ||\varphi||_D
\]

\subsection{The regularity claim} Riesz's representation theorem gives us a
sequence of functions \(f_1,f_2,\ldots\) that is Cauchy with respect to the
Dirichlet norm, such that for any \(g\in \Omega^0(S,\mathbb{R})\) we have
\(\lim_{i\to \infty}\langle f_i,g\rangle_D= \hat\rho(g)\).

We first show that the sequence \((f_i)\) converges to an \(L^2\) function on
\(S\). This is easy locally-- on any coordinate chart \(U\) we may assume
\(f_i\) has integral zero by adding a constant, so that we have
\(||f_i-f_j||_{L^2(U)}\leq C\cdot ||f_i-f_j||_D\) by \cref{cor:poincare-ineq}.
Thus the sequence is Cauchy with respect to the \(L^2\) norm which is complete
on the space of square-integrable functions.

To show the sequence is Cauchy in \(L^2(S)\), suppose \(U,V\) are open
coordinate charts with \(U\cap V\neq \emptyset\) and \((f_i)\) converges in
\(L^2(U)\). We know there are
constants \(c_1,c_2,\ldots\) such that \((f_i+c_i)\) converges
in \(L^2(V)\), and up to adding a fixed constant to each \(c_i\), we may also
assume the two limits agree on \(L^2(U\cap V)\). But this implies \(c_i\rightarrow 0\) so in fact we could have
chosen \(c_i=0\) to begin with. Since \(S\) is connected, it follows that the
sequence \((f_i)\) converges in the \(L^2\) norm locally on each chart, and
hence converges to some \(f\in L^2(S)\).

It remains to show \(f\) is smooth, which can be done locally. This is then the
content of the following result.

\begin{proposition}[(Weyl's lemma)]
    Suppose \(U\) is a bounded open set in \(\mathbb{C}\) and \(\rho\in
    \Omega^2(U)\) is a smooth \(2\)-form. If \(f\in L^2(U)\) is such that for
    any smooth compactly supported function \(g\) we have \(\int_U (\Delta f)g =
    \int_U g\rho\) then \(f\) is smooth and satisfies \(\Delta f = \rho\).

    \begin{proof}
        Omitted, see \cite{donaldsonRiemannSurfaces}. The first step is to
        reduce to \(\rho=0\) by convolutions, so that we have a weakly harmonic
        function \(f\) (i.e.\ an \(L^2\) function such that \(\Delta f=0\) in a
        weak sense). Then we want to show that weakly harmonic functions are
        smooth, which we do by finding an explicit inverse to the Laplace
        operator in terms of convolutions.
    \end{proof}
\end{proposition}

\section{The Riemann--Roch theorem}
\marginpar{Lecture 7, 21/11/23} Let \(S\) be a compact Riemann surface of genus \(g\), and let
\(D=\{p_1,\ldots,p_d\}\) be a set of distinct points in \(S\).

\begin{definition}
    We denote by \(\HH^0(D)\) the vector space of meromorphic functions on
    \(S\) with at worst a simple pole at each \(p_j\). Write
    \(\HH^0(K-D)\) for the vector space of holomorphic \(1\)-forms vanishing
    at each \(p_j\).
\end{definition}

The cohomological notation will be explained later. For now, these are complex
vector spaces keeping track of functions and \(1\)-forms with interesting
properties and the Riemann--Roch theorem will relate their dimensions. A partial
result in the form of an inequality was first proven by Riemann, and Roch
subsequently completed the work by computing the error term. This, together with
its many generalisations (the Hirzebruch--Riemann--Roch theorem, the
Grothendieck--Riemann--Roch theorem, and the Atiyah--Singer index theorem) is
one of the most fundamental results of algebraic geometry.

\begin{theorem}[(Riemann--Roch)]
    For \(S,D\) as above, we have
    \(\dim\HH^0(D)-\dim\HH^0(K-D)=d-g+1.\)
\end{theorem}

This is typically shown using algebro-geometric methods (Serre duality), but we
don't know that \(S\) has an algebraic structure yet. In fact we will use this
result to show all compact Riemann surfaces are algebraic, so our proof of the
Riemann--Roch theorem will be analytic.

\subsection{\(\mathbf{(0,1)}\)--cohomology classes control meromorphic
functions} To motivate the proof, we revisit a construction from before.
Suppose we want to construct a meromorphic function \(f\) on \(S\) with exactly
one pole, simple at \(p\in S\). We can easily find a smooth function
\(\psi\in \Omega^0(S\setminus \{p\}, \mathbb{C})\) that is meromorphic near
\(p\) with a simple pole; for instance by taking the function \(\frac{1}{z}\) in
a local chart and extending globally by multiplying with a cutoff function \(\beta\) which is \(1\) near \(p\) and vanishes outside a neighbourhood of \(p\). The problem is thus reduced to finding a smooth function \(g\in
\Omega^0(S,\mathbb{C})\) such that \(g-\psi\) is holomorphic
on \(S\setminus \{p\}\), or equivalently \(\bar\partial g= \bar\partial \psi\)
on \(S\setminus \{p\}\).

Consequently consider the smooth \((0,1)\)-form \(A=\bar\partial \psi\), which
vanishes near \(p\) (since \(\psi\) is holomorphic there) so extends by \(0\)
across \(p\) to give a global form on \(S\). We want to find \(g\in
\Omega^0(S,\mathbb{C})\) which solves \(\bar\partial g + A=0\), the obstruction
to this is precisely the class \([A]\in \HH^{0,1}(S)\) and we have a solution if
and only if \([A]=0\). But even if \([A]\neq 0\), this obstruction class is
defined (up to scaling) independently of choice of \(\psi\). Indeed if
\(\varphi\in \Omega^0(S\setminus\{p\},\mathbb{C})\) is meromorphic near \(p\)
with a simple pole, then for some \(\lambda \in \mathbb{C}\) (given by the
residue) we have that \(\varphi-\lambda \psi\) is smooth everywhere (including
near \(p\)) so \([\bar\partial \varphi]= \lambda[\bar\partial \psi]\in
\HH^{0,1}(S)\). Thus the obstruction to finding a global meromorphic function is
well-defined (up to scaling) as a \((0,1)\)-Dolbeault cohomology class. Moreover
we can represent this class by a \((0,1)\)--form \(\bar\partial \psi\) for
\(\psi\) supported on a small punctured neighbourhood of \(p\), smooth away from
\(p\).

For multiple distinct points \(p_1,\ldots,p_d\in S\), we get \((0,1)\)--forms \(A_j
= [\bar\partial \psi_j]\) supported near \(p_j\) from the above construction.
Similar arguments show we can find a meromorphic function with simple poles
precisely on the \(p_i\) provided there are \(\lambda_j\in \mathbb{C}\setminus
0\) with \(\lambda_1[A_1]+\ldots+\lambda_d[A_d]=0\) in \(\HH^{0,1}(S)\). Indeed, in
that case \(\psi=\sum_j \lambda_j\psi_j \in \Omega^0(S\setminus
\{p_1,\ldots,p_d\},\mathbb{C})\) is non-zero and meromorphic with a simple pole
near each \(p_j\), and \(\bar\partial \psi\) extends to a global \((0,1)\)--form
that is cohomologically trivial. Thus there is a smooth function \(g\) with
\(\bar\partial \psi = \bar\partial g\) on \(S\setminus \{p_1,\ldots,p_d\}\) and
\(\bar\partial g  (p_j)= 0\), so that \(g-\psi\) is the required meromorphic
function.

We can immediately use this reasoning to prove some results.

\begin{corollary}
    Suppose \(S\) has genus \(g\). Then given
    \(g+1\) points \(p_1,\ldots,p_{g+1}\) in \(S\), there is a meromorphic
    function on \(S\) with at worst simple poles, at least one pole, and smooth
    away from \(p_1,\ldots,p_{g+1}\).
    \begin{proof}
        We know \(\HH^{0,1}(S)\) has dimension \(g\), so any \(g+1\) elements of
        \(\HH^{0,1}(S)\) must be linearly dependent.
    \end{proof}
\end{corollary}

In particular, we have a classification of compact Riemann surfaces of genus
\(0\).

\begin{corollary}
    Suppose \(S\) has genus \(0\). Then \(S\) is biholomorphic to the Riemann
    sphere.
    \begin{proof}
        Indeed then \(\HH^{0,1}(S)=0\), so we can find a meromorphic function
        with a simple pole which gives the required biholomorphism.
    \end{proof}
\end{corollary}

\subsection{The residue map} The second important ingredient in proving the
Riemann--Roch theorem is the concept of residues. What is the residue of a
meromorphic function \(f\) at a point \(p\in S\)? In a local coordinate \(z\),
if \(f(z)=\sum_i a_i z^i\) then the residue is \(a_{-1}\). With respect to
another local coordinate \(w=\lambda_1 z + \lambda_2 z^2+\ldots\), the residue is
\(\lambda_1a_{-1}\).  Noting that \(\frac{\partial}{\partial z}=\lambda_1
\frac{\partial}{\partial w}\in T_p^\text{hol}S\) (the holomorphic tangent
space\footnote{The holomorphic tangent space can be defined as follows.
Recall that the tangent space \(T_pS\) at \(p\in S\) has a natural complex
structure \(J:T_pS\to T_pS\), which extends to a \(\mathbb{C}\)-linear complex
structure \(J:T_pS\otimes \mathbb{C}\to T_pS\otimes \mathbb{C}\) by writing
\(J(v+iw)=J(v)+iJ(w)\) for \(v,w\in T_pS\). Then \(J^2=-\text{Id}\) so \(J\) has
eigenvalues \(\pm i\), and the holomorphic tangent space is the
\(+i\)--eigenspace. Equivalently, one can check that this is precisely the
subspace of \(T_pS\otimes \mathbb{C}\) where all anti-linear cotangent-vectors
(i.e.\ elements of \(T_p^{0,1}S\) vanish.) In local coordinates, the holomorphic
tangent space has vectors of the form \(\frac{\partial}{\partial z}\) for
\(z=x+iy\) a local holomorphic coordinate.}), we thus see that the quantity
\(a_{-1}\frac{\partial}{\partial z}\in T_p^\text{hol}S\) is independent of
coordinates, and this is what we call the \textit{residue} of \(f\) at \(p\).

\subsection{The obstruction-class map} For similar reasons, we can also define
the obstruction-class at \(p\in S\) from the previous section as a map
\(A:T_p^\text{hol}S \to \HH^{0,1}(S)\) by first fixing a smooth cutoff function
\(\beta\), and then mapping \(\frac{\partial}{\partial z}\mapsto [\bar\partial
\beta\cdot \frac{1}{z}]\).  Indeed for a different local coordinate
\(w=\lambda_1z+\lambda_2z^2+\ldots\), we see that \(\frac{1}{z}-
\lambda_1\frac{1}{w}\) is holomorphic near \(p\) so \([\bar\partial \beta\cdot
\frac{1}{z}]=\lambda_1 [\bar\partial \beta\cdot \frac{1}{w}]\).

We examine the dual map \(A^T: \HH^{0,1}(S)^\ast \to (T_p^\text{hol}S)^\ast\).
Recall we can identify \(\HH^{1,0}(S)\) with \(\HH^{0,1}(S)^\ast\) via the
perfect pairing \((\alpha,\beta)\mapsto \int_S \alpha\wedge \beta\). Note also
that \((T_p^\text{hol}S)^\ast\) is simply the holomorphic cotangent space
\(T_p^{1,0}S\). But there is an obvious evaluation map
\(\text{ev}:\HH^{1,0}(S)\to T_p^{1,0}S\) sending a holomorphic \(1\)-form to its
value at \(p\), and we now show that \(A^T\) is essentially this.

\begin{lemma}
    Under the identifications above, \(A^T=2\pi i\cdot \textnormal{ev}\).
    \begin{proof}
        This is a local result, so choose a coordinate \(z\). Given a
        \((1,0)\)-form \(\alpha\), we see that \(A^T(\alpha)\) sends the tangent
        vector \(\frac{\partial}{\partial z}\) to \(\langle \alpha,
        A\frac{\partial }{\partial z}\rangle = \int_S \alpha \wedge
        (\bar\partial \beta \cdot \frac{1}{z}).\) But if \(\alpha\) is locally
        \(g(z)dz\), we see that the integral becomes \(\int_\mathbb{C}
        \bar\partial \beta\cdot \frac{g(z)}{z}dz\) (noting that \(\beta\) is
        supported on the coordinate chart around \(p\)). Choose a circle
        \(\gamma\) around \(p\) on which \(\beta=1\), so that Stokes' theorem
        gives us
        \[\int_\mathbb{C} \bar\partial \beta\cdot \frac{g(z)}{z}dz =
        \int_\gamma \frac{g(z)}{z}dz.\]
        By Cauchy's residue theorem, this is \(2\pi i\cdot  g(p)\) as required.
    \end{proof}
\end{lemma}

We can now prove the Riemann--Roch theorem.
All (co)-tangent spaces considered in the rest of this section will be
holomorphic (hence of \(\mathbb{C}\)-dimension one) and we will simply write
\(T_pS\) instead of \(T_p^\text{hol}S\) and likewise \(T_p^\ast S\) instead of
\(T_p^{1,0}S\).


\begin{proof}[Proof of the Riemann--Roch theorem]
    For a collection of points \(D=\{p_1,\ldots,p_d\}\subset S\), we
    have a residue map \(R: \HH^0(D)\to \bigoplus_{j=1}^d T_{p_j}S\)
    defined in the obvious way. The kernel of this map is the space of global
    holomorphic functions, which are all constant since \(S\) is compact.

    Likewise, we have a map \(A:\bigoplus_{j=1}^d T_{p_j}S\to \HH^{0,1}(S)\)
    coming from obstruction-classes. We showed that \(A(v_1,\ldots,v_d)=0\) if and
    only if \((v_1,\ldots,v_d)\) are the residues of a meromorphic function in
    \(\HH^0(D)\), i.e.\ that the kernel of \(A\) is precisely the image of
    \(R\). This gives us an exact sequence
    \[0\to \mathbb{C}\to \HH^0(D)\xrightarrow{\;R\;} \bigoplus_{j=1}^d T_{p_j}S
    \xrightarrow{\;A\;} \HH^{0,1}(S).\]

    The proof is now linear algebra-- we have \(\dim\img(A)+\dim\HH^0(D)=d+1\)
    from the exact sequence, and \(\dim\img(A)=\dim\HH^{0,1}(S)-\dim \ker
    A^T.\) Lastly, note that \(A^T:\HH^{1,0}(S)\to \bigoplus_{j=1}^d
    T_{p_j}^\ast S\) is the evaluation map, so its kernel is precisely
    \(\HH^0(K-D)\). The result follows.
\end{proof}

\section{The uniformisation theorem}

\begin{theorem}[(The uniformisation theorem)]
    Let \(S\) be a non-compact simply connected Riemann surface. Then \(S\) is
    biholomorphic to either \(\mathbb{C}\) or the upper half plane
    \(\mathbb{H}\).
\end{theorem}

Note we have already shown any simply connected compact Riemann surface must be
biholomorphic to \(\mathbb{P}^1\). Since the universal cover of any Riemann
surface is a simply connected one, we have the following classification result.

\begin{corollary}
    Any Riemann surface is biholomorphic to one of the following.
    \begin{enumerate}[(i)]
        \item The Riemann sphere \(\mathbb{P}^1\);
        \item the complex plane \(\mathbb{C}\);
        \item the punctured plane \(\mathbb{C}\setminus \{0\}\cong
            \mathbb{C}/\mathbb{Z}\);
        \item a quotient \(\mathbb{C}/\Lambda\) where \(\Lambda \subset
            \mathbb{C}\) is a lattice; or
        \item a quotient \(\mathbb{H}/\Gamma\) where \(\Gamma\subset
        \textnormal{PSL}(2,\mathbb{R})=\textnormal{Aut}\mathbb{H}\) is a
        discrete group acting freely on \(\mathbb{H}\).
    \end{enumerate}
    \begin{proof}
        Indeed, any Riemann surface is a quotient of its universal cover by an
        action of its fundamental group and the uniformisation theorem which
        surfaces can arise as universal covers. Examining their automorphism
        groups gives the result.
    \end{proof}
\end{corollary}

To prove the uniformisation theorem, we need the following analytic result which
allows us to invert Laplacians on \(S\) (this time non-compact).

\begin{theorem}
    Let \(S\) be a connected, simply connected, non-compact Riemann surface. If
    \(\rho\) is a real \(2\)-form of compact support on \(S\) with \(\int_S
    \rho = 0\), then there is a real function \(\varphi\) on \(S\) with
    \(\Delta \varphi = \rho\) such that \(\varphi \to 0\) at infinity.
\end{theorem}

Here we say \(\varphi\to 0\) at infinity in \(S\) if for all \(\epsilon > 0\)
there is a compact \(K\subset S\) such that \(|\varphi(x)|<\epsilon \) for all
\(x\notin K\). The proof is omitted, see \cite{donaldsonRiemannSurfaces}. It
is similar to inverting the Laplacian in the compact case, and the simply
connectedness provides compensation for the non-compactness.

\begin{proof}[Proof of the uniformisation theorem]
    Let \(S\) be a non-compact simply connected Riemann surface. As before, we
    construct a meromorphic function (viewed as a map to \(\mathbb{P}^1\)) and
    examine its image.

    \underbar{There is a meromorphic function \(f\)  with a simple pole.} Fix
    \(p\in S\), \(z\) a local coordinate around \(p\), and \(A=\bar\partial
    \cdot \frac{1}{z}\) a \((0,1)\)-form with \(\beta\) a cutoff function. Let
    \(\rho = \partial A\) so that \(\rho\) is a complex \(2\)-form with integral
    zero (Stokes' theorem). Thus by the theorem applied to the real and
    imaginary parts of \(\rho\) separately, we can find \(g\) with \(\Delta g =
    \rho\) such that the real and imaginary parts of \(g\) tend to \(0\) at
    infinity.

    Let \(a\) be the real part of \(A-\bar\partial g\). This is a
    \((0,1)\)-form so \(\bar\partial a=0\). Since \(\partial(A-\bar\partial
    g)=0\), we have \(da=0\). Thus since \(\HH^1(S)=0\) (by simply
    connectedness), there is a real valued \(\varphi\) satisfying
    \(d\varphi  =a\). So \(d\varphi = (A-\bar\partial g ) + (\bar A -
    \overline{\bar\partial g})\), \(A=\bar\partial g + \bar\partial
    \varphi\), and \(\bar\partial (\beta\cdot
    \frac{1}{z}-(g+\varphi))=0\). This gives us the required meromorphic
    function \(f=\beta\cdot \frac{1}{z}-g-\varphi\) with a simple pole at \(p\),
    such that the imaginary part of \(f\) tends to zero at infinity.

    \underbar{The map \(f\) is injective.} Let \(\pm \mathbb{H}\) denote the
    (open) upper and lower half-planes in \(\mathbb{C}\subset \mathbb{P}^1\),
    and let \(S_\pm=f^{-1}(\pm \mathbb{H})\) denote their preimages so that
    \(S_+\cup S_-\) is a dense open set in \(S\) (noting that \(f\) is an open
    map). Note \(S_\pm\) are non-empty, since \(p\) is a simple pole
    i.e.\ \(f\) is a local homeomorphism between neighbourhoods of \(p\in S\)
    and \(\infty \in \mathbb{P}^1\).

    Claim \(f_\pm:S_{\pm}\to \pm \mathbb{H}\) is surjective. Since we know it
    is an open map, it suffices to show it is also proper. We do this
    explicitly-- if \(B\subset \mathbb{H}\) is compact then there is an
    \(\epsilon>0\) such that \(\text{Im}(z)>\epsilon\) for all \(z\in B\). Then
    \(f^{-1}(B)=f_+^{-1}(B)\) is compact since \(\text{Im}f\) tends to zero at
    infinity. Similarly \(f_-\) is proper.

    We can now show \(f_\pm : S_\pm \to \pm\mathbb{H}\) is injective. Since
    \(f_+\) is a holomorphic map between Riemann surfaces, we know it has a
    degree which is (locally) constant on \(\mathbb{H}\). Showing injectivity is
    equivalent to showing this degree is \(1\), so suppose \(f_+\) has degree
    \(d_+\geq 2\). By surjectivity of \(f_+\), we see that for all \(n\in
    \mathbb{N}_{\geq 0}\) we can find \(x_n,\tilde{x}_n\in S_+\) with
    \(f(x_n)=f(\tilde{x}_n)=in\) so either \(x_n,\tilde{x}_n\) are distinct or
    \(f'(x_n)=0\). Since \(f\) tends to zero at infinity, the sequences
    \((x_n),(\tilde{x}_n)\) lie in a compact set \(K\) and hence converge to
    \(x,\tilde x\) respectively. But the sequence \((in)\) converges to
    \(\infty\in \mathbb{P}^1\), so \(f(x)=f(\tilde x)=\infty\) and hence
    \(x=\tilde x = p\). This contradicts \(f\) being a local homeomorphism with
    non-vanishing derivative on a neighbourhood of \(p\).

    So \(f\) maps \(S_{\pm}\) bijectively onto the open upper and lower half
    planes. This allows us to conclude \(f:S\to \mathbb{P}^1\) is injective.
    Indeed if \(x_1,x_2\in S\) are such that \(f(x_1)=f(x_2)=z\in \mathbb{P}^1\)
    then we can find disjoint discs \(D_1,D_2\) around \(x_1,x_2\) and a
    neighbourhood \(N\ni z\) such that \(N\subset f(D_1)\cap f(D_2)\). Choose
    \(z'\in N\cap \mathbb{H}\), this gives us distinct \(x_1'\in D_1\) and
    \(x_2'\in D_2\) with \(f(x_1')=f(x_2')=z'\) contradicting that \(f\) is
    injective on \(S_+\).

    \underbar{The image of \(f\) is biholomorphic to \(\mathbb{C}\) or
    \(\mathbb{H}\).} We have shown \(f:S\to \mathbb{P}^1\) is injective and
    biholomorphic onto its image, which is an open set containing
    \(\mathbb{H}\cup -\mathbb{H} \cup \{\infty\}=\mathbb{P}^1\setminus
    \mathbb{R}\). Since \(S\) is connected and simply connected, we conclude
    that \(\mathbb{P}^1\setminus f(S)\) is a compact interval \(I\) in
    \(\mathbb{R}\subset \mathbb{P}^1\). If \(I\) contains a single point, then
    \(f(S)\) is biholomorphic to \(\mathbb{C}\) and if \(I\) has nontrivial
    length, then \(f(S)\) is biholomorphic to \(\mathbb{H}\) (by writing down
    an explicit M\"obius map). This concludes the proof.
\end{proof}
This concludes the main theory of the course, the study of meromorphic functions
on Riemann surfaces.

\section{Line bundles and divisors}
\marginpar{Lecture 8, 28/11/23}
The terminology of meromorphic functions is seldom used in modern algebraic
geometry. Instead, a generalised theory of line bundles and sections is used.
These relate, for example, to the theory of embeddings into projective space.

Let \(f_1,\ldots,f_n\) be meromorphic functions on a compact Riemann surface \(S\),
and let \(\Delta\subset S\) be the (finite) set of poles outside which every
\(f_j\) is holomorphic. Thus we have a holomorphic map \(S\setminus \Delta
\to \mathbb{C}^n\), where `holomorphic' means holomorphic in each variable
separately.

\begin{exercise}
    Show this can be extended to a holomorphic map from \(S\) to
    \(\mathbb{P}^n=\mathbb{C}^n\cup \mathbb{P}^{n-1}\). Here we use a natural
    definition of holomorphic-- the standard charts on \(\mathbb{P}^n\) cover
    \(\mathbb{C}^n\) with holomorphic transition functions, and a smooth map is
    holomorphic if it is holomorphic locally.
\end{exercise}

Instead of understanding this map near \(\Delta\), the theory is clearer if we work with line bundles.

\begin{definition}
    A rank \(r\) holomorphic vector bundle over a Riemann surface \(S\) is
    given by an \(r+1\)-dimensional complex manifold \(E\) with a holomorphic
    map \(p:E\to S\) such that for each \(x\in S\), the fiber \(p^{-1}(x)=E_x\)
    has the structure of a complex vector space \(\mathbb{C}^r\) and there is
    an open set \(U\subset S\) containing \(x\) with a holomorphic and fiberwise
    linear isomorphism \(E|_U\cong U\times \mathbb{C}^n\).
\end{definition}

We think of the fibers \(E_x\) as `varying holomorphically with \(x\in S\)'. We
will be interested in line bundles, where the rank \(r\) is \(1\).

There are two further ways to think about vector bundles-- first through
transition functions. If \(\{U_\alpha\}\) is a cover of \(S\) which trivialises
\(E\), then we get two trivialisations on each intersection \(U_\alpha\cap
U_\beta\). Namely we have \(E|_{U_\alpha}\cong U_\alpha\times \mathbb{C}^n\)
and \(E|_{U_\beta}\cong U_\beta \times \mathbb{C}^n\), and on the
intersection these differ by a holomorphic endomorphism of \((U_\alpha \times
U_\beta)\times \mathbb{C}^n\) that is fiberwise linear. Thus we get a
holomorphic map \(\varphi_{\alpha\beta}:U_\alpha\cap U_\beta \to
GL(r,\mathbb{C})\) satisfying \(\varphi_{\alpha\alpha}=\text{id}\),
\(\varphi_{\alpha\beta}=\varphi_{\beta\alpha}^{-1}\), and
\(\varphi_{\alpha\gamma}=\varphi_{\alpha\beta}\circ \varphi_{\beta\gamma}\) on
triple intersections. Conversely given any cover \(\{U_\alpha\}\) and collection
of maps \(\{\varphi_{\alpha\beta}\}\) satisfying these compatibilities, we get a
vector bundle in the natural way.

Another approach to vector bundles is through sheaves. In geometry, we often
have assignments such as
\[
    \{\text{open sets }U\subset S\}\to \mathscr{O}(U)\coloneq \{\text{holomorphic functions on }U\}
\]
where the right hand side is a vector space, and where restrictions give linear
maps \(\mathscr{O}(U)\to \mathscr{O}(V)\) whenever \(V\subset U\). This is the
\textit{sheaf of holomorphic functions on \(S\)}. Similarly we can define
sheaves of continuous functions, smooth functions, differential forms etc.

\begin{definition}
    A \textit{presheaf of abelian groups} \(F\) on \(S\) is an assignment of
    abelian groups \(U\mapsto F(U)\) for every open set \(U\subset S\), with
    additive maps \(\text{res}_{V,U}:F(U)\to F(V)\) whenever \(V\subset U\)
    satisfying \(\text{res}_{U,U}=\text{id}\), \(\text{res}_{W,V}\circ
    \text{res}_{V,U}=\text{res}_{W,U}\) whenever defined. We call \(s\in F(U)\)
    a \textit{section}.

    A presheaf is a \textit{sheaf} if it satisfies two additional properties.
    Firstly, if we have a cover \(U=\bigcup U_\alpha\) and \(s\in F(U)\) is such
    that \(\text{res}_{U_\alpha,U}(s)=0\) for all \(\alpha\), then \(s=0\).
    Secondly if we have a cover \(U=\bigcup U_\alpha\) and sections
    \(s_\alpha\in F(U_\alpha)\) satisfying \(\text{res}_{U_\alpha\cap U_\beta,
    U_\alpha}(s_\alpha) = \text{res}_{U_\alpha\cap U_\beta, U_\beta}(s_\beta)\),
    then there is a section \(s\in F(U)\) satisfying
    \(\text{res}_{U_\alpha,U}(s)=s_\alpha\).
\end{definition}

\begin{example}
    A section of a vector bundle \(p:E\to S\) over \(U\subset S\) is a
    holomorphic map \(s:U\to E|_U\) satisfying \(p\circ s =\text{id}\). This
    gives a sheaf of sections of \(E\), which we denote by \(\mathscr{O}(E)\).

    The sheaf of holomorphic functions, denoted \(\mathscr{O}\), is naturally
    the sheaf of sections of the \textit{trivial line bundle}
    \(\mathbb{C}\times S\to S\).
\end{example}

\begin{example}
    One can form the tangent bundle (written \(-K\) or \(T_S\)) has fibers given
    by the holomorphic tangent spaces. Likewise, the cotangent bundle \(K\) (or
    \(T^\ast_S\)) has fibers given by the cotangent spaces, and this is the
    bundle dual to \(T_S\). Sections of \(T_S\) and \(T^\ast_S\) are holomorphic
    vector fields and \(1\)-forms respectively.
\end{example}


\subsection{The Picard group} A line bundle \(L\) on \(S\) naturally has a
\textit{dual} \(L\dual\) . This is most easily defined from the transition
functions perspective, where the transition functions for \(L\dual\) are inverse
pointwise to those of \(L\). Likewise we can take tensor products of line
bundles by multiplying transition functions. 

Since \(\mathscr{O}\) has trivial
transition functions, isomorphism classes of line bundles on \(S\) form a group
called the \textit{Picard group of \(S\)}.

We want to relate line bundles to meromorphic functions. For \(p\in S\), there
is a line bundle \(L[p]\) whose sections is given by meromorphic functions which
have at worst a simple pole at \(p\) and are holomorphic elsewhere. In terms of
transition functions, we have a trivialising cover \(S=D\cup (S\setminus
\{p\})\) where \(D\) is a disc around \(p\). Then our line bundle is determined
by a holomorphic map \(D\cap (S\setminus \{p\})\to \mathbb{C}^\ast\), i.e.\ a
nowhere vanishing function on \(D\setminus \{p\}\). We can simply choose a local
coordinate \(z\) for this, and check that changing the choice gives isomorphic
bundles.

\begin{definition}
    A divisor on a compact Riemann surface is a formal sum \(D=\sum_{j=1}^n
    a_jp_j\) for \(a_j\in \mathbb{Z}\), \(p_j\in S\). Corresponding to this we
    have the line bundle \(L[D]=\bigotimes_{j=1}^n L[p_j]^{\otimes a_j}\) where
    a negative power corresponds to taking duals first.
\end{definition}

Divisors form an abelian group \(\text{Div}(S)\) with a group homomorphism
\(\text{deg}:\text{Div}(S)\to \mathbb{Z}\) sending \(\sum_j a_j p_j\mapsto
\sum_j a_j\). Moreover, the set of divisors has a partial order such that
\(D\geq E\) if all the coefficients of \(D-E\) are non-negative.

If \(f\) is a meromorphic function with zeros of order \(a_j\) at \(p_j\) and
poles of order \(b_j\) at \(q_j\), we get a divisor \((f)=\sum a_jp_j - \sum
b_j q_j\). Divisors arising in this way are called \textit{principal}, and the
principal divisors form a subgroup of \(\text{Div}(S)\). The quotient
\(\text{Cl}(S)\) is called the \textit{divisor class group}.

These structures satisfy the following additional properties.
\begin{enumerate}
    \item Note that principal divisors have degree zero. Thus the degree map
        descends to give a group homomorphism \(\text{Cl}(S)\to \mathbb{Z}\).
    \item Two divisors \(D_1,D_2\) define isomorphic line bundles if and only if
        \(D_1-D_2\) is principal.
    \item The global holomorphic sections of \(L[D]\) correspond to meromorphic
        functions \(f\) satisfying \((f)\leq D\). Thus for \(D=\sum a_jp_j\),
        the vector space \(\HH^0(D)\) of meromorphic functions with at worst
        poles of order \(a_j\) at \(p_j\) is the space of global sections of
        \(L[D]\). We also write \(\HH^0(L[D])\) for \(\HH^0(D)\).
\end{enumerate}

The first of these implies that we have a map \(\text{Cl}(S)\to \text{Pic}(S)\).
This is in fact an isomorphism, see \cite{donaldsonRiemannSurfaces}. Thus
all line bundles arise from divisors, and we can talk of the degree of a line
bundle.


We can now restate Riemann--Roch as
\[
    \dim \HH^0(L) - \dim \HH^0(K-L) = \deg L - g + 1
\]
where \(K-L\) denotes \(T^\ast S \otimes L^\ast\). To get this form of
Riemann--Roch, we have used the identification of global sections with
meromorphic functions and a similar result that sections of \(K-L[D]\) are
holomorphic \(1\)-forms vanishing along \(D\).

\subsection{Embeddings in projective space} Let \(L\to S\) be a line bundle on a
compact Riemann surface.

\begin{definition}
    A \textit{trivialisation} of \(L\) over an open set \(U\subset S\) is a
    nowhere vanishing section of \(L\) on \(U\).
\end{definition}

Let \(s_0,\ldots,s_n\) be a basis of global sections of \(L\), and suppose for all
\(p\in S\) there is an \(s_j\) with \(s_j (p) \neq 0\) (``\(L\) is
basepoint-free''). Let \(\xi\) be a trivialisation of \(L\) over
\(U\subset S\), so \(s_j=f_j\xi\) for holomorphic functions \(f_j\) on
\(U\). Thus we have a map \(\varphi_L:U\to \mathbb{P}^n\), this is well-defined
independent of trivialisation since another trivialisation would rescale all
\(f_j(p)\) by the same amount. In particular these glue across open sets,
and we have a map \(\varphi_L: S\to \mathbb{P}^n\) written
\(\varphi_L(p)=[s_0(p):\ldots:s_n(p)]\).

\begin{definition}
    We say \(L\) is very ample if \(\varphi_L\) is a closed immersion.
\end{definition}

\begin{theorem}
    Every compact Riemann surface has a very ample line bundle. In particular,
    all compact Riemann surfaces are projective algebraic varieties.

    \begin{proof}[Sketch]
        The main idea is that once \(\text{deg}(D)\geq 2g\) then \(L[D]\) is
        basepoint free, and once \(\text{deg}(D)\geq 2g+1\) then \(L[D]\) is
        very ample.

        Take divisor \(D\) with degree \(\geq 2g+1\). Note \(\deg K = 2g-2\),
        since we showed that a holomorphic \(1\)-form has \(\chi(S)=2g-2\)
        zeros with multiplicity. Thus \(\deg(K-D)\leq -1\) and there are no
        global holomorphic sections of \(K-D\), giving us
        \(\HH^0(K-L[D])=0\). Thus by Riemann--Roch we have
        \(\dim\HH^0(L[D])\geq g+2\). Similarly for any point \(p\in S\), we can
        show \(\dim \HH^0(L[D-p])< \dim \HH^0(L[D])\) i.e.\ not all sections of
        \(L[D]\) vanish at \(p\).  Thus \(L[D]\) is basepoint-free.

        The proof of very-ampleness is analogous. It is then a very general fact
        (Chow's theorem) that all compact complex submanifolds of
        \(\mathbb{P}^n\) admit the structure of projective varieties. For
        Riemann surfaces, the proof goes roughly as follows-- if \(S\subset
        \mathbb{P}^n\) is one dimensional then the homogeneous coordinates
        \(z_i\) induce meromorphic functions  \(f_j=\frac{z_j}{z_0}\) on \(S\).
        These meromorphic functions realise \(S\) as a branched cover
        of \(\mathbb{P}^1\), so that the field of meromorphic functions on
        \(S\) is a finite field extension of the corresponding field on
        \(\mathbb{P}^1\). In particular any two meromorphic functions on
        \(S\) are related by a polynomial, so \(S\) is contained in the
        vanishing set of some polynomials. A little topology shows \(S\) is
        genuinely the vanishing locus of these polynomials.
    \end{proof}
\end{theorem}

This marks the end of the course. With another lecture we would have looked at
sheaf cohomology using Cech covers and explained Riemann--Roch in these terms,
writing \[\chi(L)\coloneq \dim \HH^0(L)-\dim\HH^1(L)=\deg L - g + 1.\] This is
proven through Serre duality, \(\HH^0(K-L)\cong \HH^1(L)\). The motivation for
cohomology is explained in \cite[chapter 12]{donaldsonRiemannSurfaces}, and
this clarifies many constructions in the course --- in particular how to view
higher cohomology groups as obstructions.
